% concatenated from main.tex
\documentclass[12pt,leqno,twoside,sort]{amsart}

\usepackage{amssymb,amsmath,amsthm,soul,color}

\usepackage[normalem]{ulem}

\usepackage{amssymb,amsmath,amsthm}

\usepackage{mathrsfs}

\usepackage[utf8]{inputenc}

\usepackage{todonotes}
\setuptodonotes{inline}

\usepackage[left=2cm, right=2cm, top=2cm, bottom=2cm]{geometry}

\usepackage{url}

\usepackage[numbers,sort&compress]{natbib}

\usepackage{fixmath}

\usepackage[T1]{fontenc}
\usepackage{lmodern}
\usepackage{microtype}
\usepackage[normalem]{ulem}
\usepackage{mathtools}

\usepackage[colorlinks=true,urlcolor=blue,
citecolor=red,linkcolor=blue,linktocpage,pdfpagelabels,
bookmarksnumbered,bookmarksopen]{hyperref}


\usepackage{bbm}

\usepackage{xcolor,cancel}

\newcommand\hcancel[2][black]{\setbox0=\hbox{$#2$}%
\rlap{\raisebox{.45\ht0}{\textcolor{#1}{\rule{\wd0}{1pt}}}}#2}

\linespread{1.2}

\newtheorem{Th}{Theorem}[section]
\newtheorem{Prop}[Th]{Proposition}
\newtheorem{Lem}[Th]{Lemma}
\newtheorem{Cor}[Th]{Corollary}
\newtheorem{Def}[Th]{Definition}
\newtheorem{Rem}[Th]{Remark}
\newtheorem{Ex}[Th]{Example}
\newtheorem{Ass}[Th]{Assumption}
\newtheorem{Que}[Th]{Question}


\newenvironment{altproof}[1]
{\noindent
{\em Proof of {#1}}.}
{\nopagebreak\mbox{}\hfill $\Box$\par\addvspace{0.5cm}}

\newcommand{\wt}{\widetilde}
\newcommand{\mf}{\mathfrak}
\newcommand{\vp}{\varphi}
\newcommand{\Vp}{\Phi}
\newcommand{\eps}{\varepsilon}
\newcommand{\ov}{\overline}

\newcommand{\R}{\mathbb{R}}
\newcommand{\Q}{\mathbb{Q}}
\newcommand{\Z}{\mathbb{Z}}
\newcommand{\N}{\mathbb{N}}

\newcommand{\F}{{\mathcal F}}

\newcommand{\cA}{{\mathcal A}}
\newcommand{\cB}{{\mathcal B}}
\newcommand{\cC}{{\mathcal C}}
\newcommand{\cD}{{\mathcal D}}
\newcommand{\cE}{{\mathcal E}}
\newcommand{\cF}{{\mathcal F}}
\newcommand{\cG}{{\mathcal G}}
\newcommand{\cH}{{\mathcal H}}
\newcommand{\cI}{{\mathcal I}}
\newcommand{\cJ}{{\mathcal J}}
\newcommand{\cK}{{\mathcal K}}
\newcommand{\cL}{{\mathcal L}}
\newcommand{\cM}{{\mathcal M}}
\newcommand{\cN}{{\mathcal N}}
\newcommand{\cO}{{\mathcal O}}
\newcommand{\cP}{{\mathcal P}}
\newcommand{\cQ}{{\mathcal Q}}
\newcommand{\cR}{{\mathcal R}}
\newcommand{\cS}{{\mathcal S}}
\newcommand{\cT}{{\mathcal T}}
\newcommand{\cU}{{\mathcal U}}
\newcommand{\cV}{{\mathcal V}}
\newcommand{\cW}{{\mathcal W}}
\newcommand{\cX}{{\mathcal X}}
\newcommand{\cY}{{\mathcal Y}}
\newcommand{\cZ}{{\mathcal Z}}
\newcommand{\al}{\alpha}
\newcommand{\be}{\beta}
\newcommand{\ga}{\gamma}
\newcommand{\de}{\delta}
\newcommand{\De}{\Delta}
\newcommand{\Ga}{\Gamma}
\newcommand{\Om}{\Omega}

\newcommand{\weakto}{\rightharpoonup}
\newcommand{\pa}{\partial}

\newcommand{\curl}{\nabla \times}
\renewcommand{\div}{\mathrm{div}\,}
\newcommand{\sgn}{\mathrm{sgn}}
\newcommand{\spann}{\mathrm{span}\,}
\newcommand{\range}{\mathrm{range}\,}

\newcommand{\tX}{\widetilde{X}}
\newcommand{\tu}{\widetilde{u}}
\newcommand{\tv}{\widetilde{v}}
\newcommand{\tcV}{\widetilde{\cV}}

\newcommand{\wh}{\widehat}
\newcommand{\wti}{\widetilde}
\newcommand{\cTto}{\stackrel{\cT}{\longrightarrow}}
\newcommand{\ctto}{\stackrel{\tau}{\longrightarrow}}

\numberwithin{equation}{section}

\DeclareMathOperator*{\esssup}{ess\,sup}
\DeclareMathOperator*{\essinf}{ess\,inf}

\DeclareMathOperator*{\argmin}{arg\,min}

\newcommand{\supp}{\mathrm{supp}\,}

\newcommand{\loc}{\mathrm{loc}}

\newcommand{\vertiii}[1]{{\left\vert\kern-0.25ex\left\vert\kern-0.25ex\left\vert #1 
    \right\vert\kern-0.25ex\right\vert\kern-0.25ex\right\vert}}

\begin{document}

\title[Nonlinear scalar field equations with a critical Hardy potential]{Nonlinear scalar field equations with a critical Hardy potential}

\author[B. Bieganowski]{Bartosz Bieganowski}
\author[D. Strzelecki]{Daniel Strzelecki}

\address[B. Bieganowski]{\newline\indent
	Faculty of Mathematics, Informatics and Mechanics, \newline\indent
	University of Warsaw, \newline\indent
	ul. Banacha 2, 02-097 Warsaw, Poland}
\email{\href{mailto:bartoszb@mimuw.edu.pl}{bartoszb@mimuw.edu.pl}}

\address[D. Strzelecki]{\newline\indent
	Faculty of Mathematics, Informatics and Mechanics, \newline\indent
	University of Warsaw, \newline\indent
	ul. Banacha 2, 02-097 Warsaw, Poland}
\email{\href{mailto:dstrzelecki@mimuw.edu.pl}{dstrzelecki@mimuw.edu.pl}}

\date{}	\date{\today} 

\begin{abstract} 
We study the existence of solutions for the nonlinear scalar field equation
$$-\Delta u - \frac{(N-2)^2}{4|x|^2} u  = g(u),  \quad \mbox{in } \R^N \setminus \{0\},$$
where the potential $-\frac{(N-2)^2}{4|x|^2}$ is the critical Hardy potential and $N \geq 3$. The nonlinearity $g$ is continuous and satisfies general subcritical growth assumptions of the Berestycki-Lions type. The problem is approached using variational methods within a non-standard functional setting. The natural energy functional associated with the equation is defined on the space $X^1(\R^N)$, which is the completion of $H^1(\R^N)$ with respect to the norm induced by the quadratic part of the functional. We establish the existence of a nontrivial solution $u_0 \in X^1(\R^N)$ that satisfies the Poho\v{z}aev constraint $\cM$ and minimizes the energy functional on $\cM$. Furthermore, assuming $g$ is odd, we prove the existence of at least one non-radial solution.

\medskip

\noindent \textbf{Keywords:} variational methods, singular potential, scalar field equation, critical Hardy potential
   
\noindent \textbf{AMS Subject Classification:} 35Q99, 35J10, 35J20, 58E99
\end{abstract}

\maketitle

\section{Introduction}

We study the existence, multiplicity, and qualitative properties of solutions for the nonlinear scalar field equation
\begin{equation}\label{eq:intro_main}
-\Delta u - \mu \frac{u}{|x|^2} = g(u), \quad \mbox{in } \R^N \setminus \{0\},
\end{equation}
where $N \geq 3$. This paper focuses specifically on the \textit{critical case} $\mu = \mu_c = \frac{(N-2)^2}{4}$, where $\mu_c$ is the best constant in the classical Hardy inequality. The inverse-square potential $\frac{1}{|x|^2}$, commonly referred to as the Hardy potential, is of significant importance in quantum mechanics, modeling diverse phenomena such as the interaction of a charge with a dipole or certain molecular potentials \cite{Alhaidari2014, Bawin, Case}. The nonlinearity $g$ is assumed to satisfy the general Berestycki-Lions conditions, (G1)--(G5) below, which do not include the celebrated Ambrosetti-Rabinowitz condition.

This problem presents a combination of two significant analytical challenges: the \textit{criticality of the potential} and the \textit{generality of the nonlinearity}.

The foundational case, where $\mu = 0$, was famously studied by Berestycki and Lions in their seminal work \cite{BL}. They established the existence of a ground state solution in the standard Sobolev space $H^1(\R^N)$ under these general conditions. The inclusion of the \textit{critical} Hardy potential $\mu_c / |x|^2$, however, fundamentally alters the problem's functional analytic structure. The associated quadratic form $\xi(u) = \int_{\R^N} (|\nabla u|^2 - \mu_c |x|^{-2} u^2) \, dx$ is merely positive semi-definite, not coercive, on $H^1(\R^N)$. This loss of coercivity necessitates moving the analysis to a larger, Sobolev-type space, $X^1(\R^N)$, which is the completion of $H^1(\R^N)$ with respect to the problem's natural norm. The properties of this space, which is essential for our analysis, build on foundational work by authors such as Frank \cite{Frank}, Suzuki \cite{Suzuki}, Mukherjee, Nam, and Nguyen \cite{Mukherjee}, Trachanas and Zographopoulos \cite{TrachanasZographopoulos, TrachanasZographopoulosCor}.

Our focus on the \textit{critical} case $\mu = \mu_c$ should be contrasted with the \textit{subcritical} case $\mu < \mu_c$. For subcritical potentials, the quadratic form remains coercive (or can be made so with an equivalent norm), allowing for a more standard application of variational methods within $H^1(\R^N)$ or related weighted spaces. The subcritical case was of interest to many authors (see e.g. \cite{GuoMederski, LiLiTang, MR3989664} and references therein). For instance, Li, Li, and Tang \cite{LiLiTang} studied the subcritical problem under Berestycki-Lions-type conditions, establishing the existence of a ground state solution \textit{within} the $H^1(\R^N)$ space. We find that for the critical potential, any nonnegative solution is genuinely singular and \textit{does not} belong to $H^1(\R^N)$. This highlights a fundamental difference in the qualitative nature of solutions as the potential transitions from subcritical to critical. Other related works, such as those by Deng, Jin, and Peng \cite{Deng}, Felli, Pistoia \cite{Felli}, Smets \cite{Smets}, and Terracini \cite{Terracini96}, have investigated problems combining the Hardy potential with the \textit{critical Sobolev exponent} $2^* = \frac{2N}{N-2}$, which introduces a different set of compactness challenges related to the nonlinearity itself.

The second major challenge, the absence of the Ambrosetti-Rabinowitz condition, means that standard mountain pass arguments do not guarantee the boundedness of Palais-Smale sequences. For the classical problem ($\mu=0$), this difficulty is often overcome by seeking solutions on the Pohožaev constraint. A systematic approach to find such solutions was developed by Jeanjean \cite{Jeanjean}, and later refined by Jeanjean, Tanaka \cite{JeanjeanTanaka} and Hirata, Ikoma, Tanaka \cite{HirataIkomaTanaka}. They introduced an \textit{augmented functional} technique to construct a bounded Poho\v{z}aev-Palais-Smale sequence at the mountain pass level. We adopt this strategy to construct a bounded Poho\v{z}aev-Palais-Smale sequence for our singular problem.


The mere boundedness of a (PPS) sequence is not sufficient; one must additionally establish its convergence. The primary tool for this in translation-invariant problems is the profile decomposition method, developed by Gérard \cite{Gerard} and Nawa \cite{Nawa}. This method fails in our setting because the $X^1(\R^N)$ space, due to the potential centered at the origin, is \textit{not} translation-invariant. A central technical contribution of this paper is the development of a \textit{new} profile decomposition specifically adapted to bounded (PPS) sequences in this non-translation-invariant space. This decomposition is crucial as it shows that any potential \textit{bubbles} (weak limits of translated sequences) must be solutions to the limiting problem \textit{without} the potential, i.e., in the classical $H^1(\R^N)$ space. By showing that our mountain pass level is strictly less than the ground state energy of this limiting problem (an energy comparison adapting \cite{JeanjeanTanaka}), we prove that no bubbles can form, forcing the strong convergence of the sequence to our solution.

Finally, we extend our analysis to the existence of non-radial solutions. This is a classical and challenging question, famously posed by Berestycki and Lions \cite{BL} for the $\mu=0$ case. This question was recently answered for the classical problem by Mederski \cite{MR4173560}, who found non-radial solutions by minimizing the functional on a subspace of $H^1(\R^N)$-functions orthogonal to the space of radial functions. We successfully adapt this symmetry-based minimization argument of \cite{MR4173560} from the standard $H^1(\R^N)$ space to our non-translation-invariant, singular $X^1(\R^N)$ setting, thereby establishing the existence of non-radial solutions for the critical Hardy potential problem under general nonlinearities.


In the critical case, \eqref{eq:intro_main} rewrites as
\begin{equation}\label{eq:main}
-\Delta u - \frac{(N-2)^2}{4|x|^2} u  = g(u),  \quad \mbox{in } \R^N \setminus \{0\}.
\end{equation}
We consider the following general conditions
\begin{itemize}
\item[(G1)] $g : \R \rightarrow \R$ is continuous;
\item[(G2)] $-\infty < \liminf_{s \to 0}  \frac{g(s)}{s} \leq \limsup_{s \to 0} \frac{g(s)}{s} = -\ell < 0$;
\item[(G3)] $\lim_{|s| \to \infty} \frac{g(s)}{|s|^{p-1}} = 0$ and $\lim_{|s| \to \infty} g(s)s \geq 0$ for some $p \in \left( 2,2^* \right)$;
\item[(G4)] there is $\zeta_0 > 0$ such that $G(\zeta_0) > 0$, where $G(s) := \int_0^{s} g(\tau) \, d\tau > 0$;
\item[(G5)] there are $\nu > 2$ and $\gamma_0 < 0$ such that
$G(s) - \frac{1}{\nu} g(s)s \leq \gamma_0 s^2$ for $s\in\R$.
\end{itemize}

Assumptions (G1)--(G4) are standard, in the spirit of \cite{BL}, although we do not require oddness of $g$. The condition (G5) is new and is required to verify the boundedness of (PPS) sequences. It can be treated as a variant of the Ambrosetti-Rabinowitz condition, which in its original form
\begin{equation}\tag{AR}\label{AR}
\mbox{there is } \nu > 2 \mbox{ such that } 0 \leq G(s) \leq \frac{1}{\nu} g(s)s, \quad s \in \R
\end{equation}
is not necessarily satisfied under (G1)--(G5).

We consider the so-called \textit{positive-mass case} ($\ell > 0$). Therefore, it is convenient to define $f : \R \rightarrow \R$ by
\begin{equation}\label{def:f}
f(s) := \ell s + g(s), \quad s \in \R
\end{equation}
and then $G$ can be rewritten as
$$
G(s) = -\frac{\ell}{2} s^2 + F(s), \quad F(s) := \int_0^s f(\tau) \, d\tau.
$$
Then conditions (G1)--(G4) imply the following
\begin{itemize}
\item[(F1)] $f : \R \rightarrow \R$ is continuous;
\item[(F2)] $-\infty < \liminf_{s \to 0}  \frac{f(s)}{s} \leq \limsup_{s \to 0} \frac{f(s)}{s} = 0 < \ell$;
\item[(F3)] $\lim_{|s| \to \infty} \frac{f(s)}{|s|^{p-1}} = 0$ and $\lim_{|s| \to \infty} f(s)s \geq 0$ for some $p \in \left( 2, 2^* \right)$;
\item[(F4)] there is $\zeta_0 > 0$ such that $F(\zeta_0) > \frac{\ell}{2} \zeta_0^2$;
\item [(F5)] there are $\nu > 2$ and $\gamma_0 < \ell\left(\frac 12-\frac 1\nu\right)$ such that
$F(s) - \frac{1}{\nu} f(s)s \leq \gamma_0 s^2$ for $s\in\R$.
\end{itemize}
Observe that $f$ also doesn't need to satisfy the \eqref{AR} condition. In fact, \eqref{AR} implies (F5). Below we present a few examples of nonlinearities satisfying (F1)--(F5).

\begin{Ex}
In what follows, we fix $p \in (2,2^*)$.
\begin{enumerate}
    \item[(a)] It is standard to check that $f(s)= |s|^{q-2}s$ satisfies (F1)--(F4) with $2 < q < p$ in (F3). To see (F5) note that
    $$
    F(s) - \frac{1}{q} f(s)s = \frac{1}{q} |u|^q - \frac1q |u|^q = 0. 
    $$
    \item[(b)] Consider $f(s) = \ln (e + |s|) |s|^{q-2}s$, where $2 < q < p$. It is clear that (F1)--(F3) are satisfied. Since $f$ is bounded from below by the power nonlinearity, (F4) is obvious. To get (F5) observe that, for $\nu=q$
    \[
    F(s)-\frac{1}q f(s)s\leq \ln(e+|s|)\left(\int_0^s |s|^{q-2}s\,ds-\frac1q |s|^q \right)=0.
    \]
   
    \item[(c)] Let $\varepsilon > 0$, $2 < q < p$, $0 < \alpha < q-2$. Consider $f(s) = \frac{|s|^{q-2}s}{1+\varepsilon |s|^\alpha}$. Clearly (F1)--(F3) are satisfied. To see (F4), we compute for sufficiently large $s > 1$,
    $$
    F(s) = \int_0^s \frac{|\tau|^{q-2} \tau}{1+\varepsilon |\tau|^\alpha} \, d\tau \geq \frac{1}{1+\varepsilon s^\alpha} \int_0^s \tau^{q-1} \,d\tau = \frac{s^q}{q (1+\varepsilon s^\alpha)} > \frac{s^{q-\alpha}}{q (1+\varepsilon )} = \frac{s^{q-\alpha-2}}{q (1+\varepsilon )} s^2 > \frac{\ell}{2} s^2.
    $$
    To get (F5) we note that $s \mapsto H(s) := F(s) - \frac{1}{\nu} f(s)s$ is even, so it is enough to verify (F5) for $s > 0$. For such $s$ we compute
    \begin{align*}
    H'(s) &= \frac{s^{q-1}}{1+\varepsilon s^\alpha} - \frac{1}{\nu} \frac{q s^{q-1} (1+\varepsilon s^\alpha) - \alpha \varepsilon s^{q+\alpha - 1}}{(1+\varepsilon s^\alpha)^2} = \frac{1}{\nu} \frac{s^{q-1}}{1+\varepsilon s^\alpha} \left( \nu - q + \frac{ \alpha \varepsilon s^\alpha }{1+\varepsilon s^\alpha} \right) \\
    &< \frac{1}{\nu} \frac{s^{q-1}}{1+\varepsilon s^\alpha} \left( \nu - q + \alpha \right) = 0
    \end{align*}
    for $\nu = q - \alpha > 2$. Taking into account that $H(0) = 0$, we obtain
    $$
    F(s) - \frac{1}{q-\alpha} f(s)s \leq H(0) = 0.
    $$
    
    \item[(d)] Consider $f(s) = (2 + \arctan(s)) |s|^{q-2}s$ with $2 < q < p$. Since $s \mapsto 2+\arctan s$ is increasing and bounded away from zero, it satisfies $(F1)-(F5)$ in the same way as in (b).
    \item[(e)] Consider $2 < r < q < p$ and
    $$
    f(s) = |s|^{q-2}s - |s|^{r-2}s.
    $$
    Clearly (F1)--(F4) are satisfied. To see (F5) note that
    $$
    F(s) - \frac{1}{q} f(s)s = \frac1q |s|^q - \frac1r |s|^r - \frac1q |s|^q + \frac1q |s|^r = \left( \frac1q - \frac1r \right) |s|^r < 0
    $$
    for $s \neq 0$. Note that $f$ doesn't satisfy \eqref{AR}.
    \item[(f)] Let $2 < q < p$ and
    $$
    f(s) = |s|^{q-2}s + e^{-1/|s|} \frac{|s|^{q-2}s}{1 + |s|^{q-2}}.
    $$
    It is easy to check that (F1)--(F4) hold. To show (F5), as in (c) we have that $s \mapsto H(s) := F(s) - \frac{1}{q} f(s)s$ is even. Hence, let $s > 0$. Note that the power part $|s|^{q-2}s$ of $f$ does not give any contribution in $H$, so 
    $$
    H(s) = \int_0^s e^{-1/\tau} \frac{\tau^{q-1}}{1 + \tau^{q-2}} \, d\tau - \frac{1}{q} e^{-1/s} \frac{s^{q}}{1 + s^{q-2}}.
    $$
    Then
    \begin{align*}
    H'(s) &= e^{-1/s} \frac{s^{q-1}}{1 + s^{q-2}} - \frac{1}{q} e^{-1/s} \frac{2s^{2q+1} + s^{2q} + q s^{q+3} + x^{q+2}}{(s^q + s^2)^2} \\ &= \frac{e^{-1/s} \left( (q-2) s^{2q+1} - s^{2q} - s^{q+2} \right)}{q (s^q+s^2)^2} \leq \frac{(q-2) s^{2q+1}}{q (s^q+s^2)^2} \leq \frac{q-2}{q} s.
    \end{align*}
    Therefore
    $$
    H(s) \leq \frac{q-2}{2q} s^2 =: \gamma_0 s^2.
    $$
    Hence the assumptions are satisfied with
    $$
    \ell > \gamma_0 \left(\frac12 - \frac1q \right)^{-1} = 1.
    $$
Note that this example does not satisfy \eqref{AR}. Indeed, for $H_\nu(s):=F(s)-\frac{1}{\nu}f(s)$ we have
\[
H_\nu'(s)=\frac{e^{-1/s}\left(((\nu-2)s - 1)s^{2q}+s^{q+2}(\nu s - qs - 1) \right)}
     {\nu(s^{q} + s^{2})^{2}}.
\]
Hence, for $\nu>2$ there is $\lim_{s\to \infty} H_\nu'(s)=\infty$, so $\lim_{s\to\infty} H_\nu(s)=\infty$ and \eqref{AR} is not satisfied for any $\nu > 2$.
\end{enumerate}
\end{Ex}

The natural energy functional associated with \eqref{eq:main} is given by
\begin{align*}
J (u) &= \frac{1}{2} \int_{\R^N} |\nabla u|^2 - \frac{(N-2)^2}{4} \frac{u^2}{|x|^2} \, dx - \int_{\R^N} G(u) \, dx \\
&=  \frac{1}{2} \int_{\R^N} |\nabla u|^2 - \frac{(N-2)^2}{4} \frac{u^2}{|x|^2} \, dx + \frac{\ell}{2} \int_{\R^N} u^2 \, dx - \int_{\R^N} F(u) \, dx.
\end{align*}
Clearly, $J$ is well defined on $H^1 (\R^N)$, however the quadratic form
$$
u \mapsto \|u\|^2 := \xi(u) +  \int_{\R^N} u^2 \, dx,
$$
where
$$
\xi(u) := \int_{\R^N} |\nabla u|^2 - \frac{(N-2)^2}{4} \frac{u^2}{|x|^2} \, dx,
$$
generates a norm in $H^1 (\R^N)$, which is not complete (and, in particular, is not equivalent to the standard one). Therefore, to work in a Banach space, we need to slightly enlarge the energy space. More precisely, we define
$$
X^1 (\R^N) := \mbox{completion of } H^1(\R^N) \mbox{ with respect to the norm } \| \cdot \|.
$$
For properties of $X^1 (\R^N)$ and the variational setting (including the appropriate definition of $\xi$ and $\|\cdot\|$), see Sections \ref{sect:2} and \ref{sect:3}, respectively. Under the foregoing assumptions, $J$ is a $\cC^1$ functional on $X^1(\R^N)$ and its critical points are weak solutions to \eqref{eq:main}.

It is natural to consider the Poho\v{z}aev constraint associated with \eqref{eq:main}, namely
$$
\cM := \left\{ u \in X^1 (\R^N) \setminus \{0\} \ : \ \xi(u) = 2^* \int_{\R^N} G(u) \, dx \right\}.
$$
In the classical setting it is known that every weak solution of a scalar field equation belongs to the Poho\v{z}aev constraint. Here, due to the presence of the critical Hardy potential, it is unclear whether all nontrivial critical points of $J$ belong to $\cM$. In our analysis we will find a solution that minimizes $J$ on $\cM$, although it is vague whether it is a ground state solution.

We are ready to state our first result.

\begin{Th}\label{Th:Main1-Existence}
Suppose that (G1)--(G5) are satisfied. 
\begin{enumerate}
\item[(a)] There exists a nontrivial solution $u_0 \in X^1 (\R^N)$ of \eqref{eq:main} that satisfies the Poho\v{z}aev constraint $\cM$ and $J(u_0) > 0$. 
\item[(b)] $u_0$ is a minimizer of the energy functional $J$ on $\cM$.
\item[(c)] Every solution $u\in X^1(\R^N)$ of \eqref{eq:main} belongs to $C^{1,\alpha}_{\loc} (\R^N \setminus \{0\})$ for some $\alpha \in (0,1)$.
\item[(d)] If, in addition, $g$ is H\"older continuous, every solution $u\in X^1(\R^N)$ of \eqref{eq:main} is of the class $C^2 (\R^N\setminus \{0\})$.
\item[(e)] If, in addition, $g$ is odd, there exists a nontrivial, nonnegative, radial and radially nonincreasing solution in $X^1 (\R^N) \cap C^2 (\R^N \setminus \{0\})$ of \eqref{eq:main}, which is a minimizer of $J$ on $\cM$.
\item[(f)] Every nonnegative solution of \eqref{eq:main} in $X^1(\R^N) \cap C^2 (\R^N \setminus\{0\})$ does not belong to $H^1 (\R^N)$.
\end{enumerate}
\end{Th}

We would like to point out that if $g$ is odd, even if we restrict the setting to radially symmetric functions, it is unclear whether a direct minimization on $\cM$ can be applied. Specifically, it is unclear whether a minimizing sequence on $\cM$ (that is, sequence $(u_n) \subset \cM$ such that $J(u_n) \to \inf_{\cM} J$) is bounded.

In this paper, we present a profile decomposition approach, which is based on the concentration-compactness principle in the spirit of Lions. The decomposition of minimizing sequences, in the spirit of \cite{Nawa, Gerard} is a challenging problem, as the space $X^1 (\R^N)$ is not translation-invariant, and even if the sequence has a weak limit after finding proper translations, it should be a solution to the limiting problem without the singular potential, for which the proper energy space is $H^1 (\R^N)$, not $X^1 (\R^N)$.

Next, we investigate the existence of non-radial solutions. We follow the ideas of \cite{MR4173560}. Suppose that $N \geq 4$ and write $\R^N = \R^M \times \R^M \times \R^{N-2M}$, where $2 \leq M \leq \frac{N}{2}$. For $x \in \R^N$ we will write $x=(x_1, x_2, x_3)$, where $x_1, x_2 \in \R^M$ and $x_3 \in \R^{N-2M}$. We introduce an action $\tau : \R^N \rightarrow \R^N$ by
$$
\tau(x_1, x_2, x_3) := (x_2, x_1, x_3)
$$
and the subspace
\begin{equation}\label{eq:Xtau}
X_\tau := \left\{ u \in X^1 (\R^N) \ : \ u(x) = - u(\tau x) \mbox{ for all } x \in \R^N \right\}.
\end{equation}
Then, if $u \in X_\tau$ is radially symmetric, then $u \equiv 0$ and therefore $X_\tau$ does not contain nontrivial, radial functions.

Let $\mathcal{O} := \mathcal{O}(M) \times \mathcal{O}(M) \times \mathrm{id}_{N-2M} \subset \mathcal{O}(N)$. We introduce the subspace $X^1_\mathcal{O} (\R^N)$ of $\mathcal{O}$-invariant functions.

\begin{Th}\label{Th:Main2}
Suppose that $N\geq 4$, (G1)--(G5) are satisfied and that $g$ is odd. Then
$$
\inf_{\cM \cap X_\tau \cap X^1_{\mathcal{O}} (\R^N)} J \geq 2 \inf_{\cM} J > 0
$$
and there exists a nontrivial solution $v_0 \in \cM \cap X_\tau \cap X^1_{\mathcal{O}} (\R^N)$ of \eqref{eq:main}, which minimizes $J$ on $\cM \cap X_\tau \cap X^1_{\mathcal{O}} (\R^N)$. In particular, $v_0$ is not radial.
\end{Th}

The paper is organized as follows. Sections \ref{sect:2} and \ref{sect:3} are devoted to the description of the functional and variational setting, respectively. The profile decomposition is proved in Section \ref{sect:4}. Then, the point (a) of Theorem \ref{Th:Main1-Existence}, namely the existence of a nontrivial solution is shown in \ref{sect:5}. Section \ref{sect:6} is devoted to proving additional properties of the solution, namely points (b) and (c). Section \ref{sect:7} contains the proof of Theorem \ref{Th:Main2}. 



In what follows, $\lesssim$ denotes an inequality up to a positive, multiplicative constant and $C$ denotes a generic, positive constant that may vary from one line to another. 




\section{Functional setting}\label{sect:2}

Recall the Hardy inequality
\begin{equation} \label{ineq:Hardy}
 \int_{\R^N} \frac{u^2}{|x|^2} \, dx \leq \frac{4}{(N-2)^2} \int_{\R^N} |\nabla u|^2 \, dx, \quad u \in H^1 (\R^N).   
\end{equation}
In particular, it implies that
$$
\| u \|^2 := \int_{\R^N} |\nabla u|^2 - \frac{(N-2)^2}{4} \frac{u^2}{|x|^2} + u^2 \, dx, \quad u \in H^1 (\R^N)
$$
is a norm in $H^1 (\R^N)$. However, due to the optimality of the constant $\frac{(N-2)^2}{4}$ and the fact that optimizers of \eqref{ineq:Hardy} do not belong to $H^1(\R^N)$\footnote{One can construct a function in $X^1 (\R^N)$, which behaves like $|x|^{- \frac{N-2}{2}}$ at the origin. Note that $x \mapsto |x|^{- \frac{N-2}{2}}$ is the optimizer of the Hardy inequality.}, this norm is not complete in $H^1 (\R^N)$. To work in a Hilbert space we will work in the completion of $H^1 (\R^N)$ in this norm. Such a space was already considered in \cite{Suzuki, Mukherjee, TrachanasZographopoulos, TrachanasZographopoulosCor}.  

Let $X^1 (\R^N)$ denote the completion of $H^1 (\R^N)$ with respect to the norm
\begin{equation}\label{normX1-oryg}
\| u \|^2 = \int_{\R^N} |\nabla u|^2 - \frac{(N-2)^2}{4} \frac{u^2}{|x|^2} + u^2\, dx, \quad u \in H^1 (\R^N).
\end{equation}
Then, from \cite[Theorem 1.2]{Frank} (see also \cite{Suzuki}), the following embeddings are continuous
$$
H^1 (\R^N) \subset X^1 (\R^N) \subset H^s(\R^N)
$$
for any $s \in [0, 1)$, where we set $H^0 (\R^N) = L^2 (\R^N)$. In particular, we have continuous embeddings into Lebesgue spaces
\begin{equation}\label{eq:embeddings}
X^1 (\R^N) \subset L^t (\R^N)
\end{equation}
for $t \in [2, 2^*)$, which are locally compact (the embedding $X^1 (\R^N) \subset L^t_{\loc} (\R^N)$ is compact).

Moreover, see \cite{TrachanasZographopoulos, TrachanasZographopoulosCor, VazquezZ}, the norm in $X^1 (\mathbb{R}^N)$ has the following characterization
\begin{align}\label{normX1}
\|u\|^2 = \lim_{\varepsilon \to 0+} \left( \int_{|x| > \varepsilon} |\nabla u|^2 \, dx - \frac{(N-2)^2}{4} \int_{|x| > \varepsilon} \frac{u^2}{|x|^2} \, dx - \frac{N-2}{2} \varepsilon^{-1} \int_{|x|=\varepsilon} u^2 \, dS \right) + \int_{\R^N} u^2 \, dx
\end{align}
for $u \in X^1 (\mathbb{R}^N)$, and the scalar product $\langle \cdot, \cdot \rangle$ on $X^1 (\R^N)$ is given by
\begin{align*}
\langle u,v \rangle = \lim_{\varepsilon \to 0+} \left( \int_{|x| > \varepsilon} \nabla u \nabla v \, dx - \frac{(N-2)^2}{4} \int_{|x| > \varepsilon} \frac{u v}{|x|^2} \, dx - \frac{N-2}{2} \varepsilon^{-1} \int_{|x|=\varepsilon} u v \, dS \right) + \int_{\R^N} uv \, dx.
\end{align*}
In what follows we introduce the notation
\begin{equation}\label{xi:uv}
\xi(u, v) := \lim_{\varepsilon \to 0+} \left( \int_{|x| > \varepsilon} \nabla u \nabla v \, dx - \frac{(N-2)^2}{4} \int_{|x| > \varepsilon} \frac{u v}{|x|^2} \, dx - \frac{N-2}{2} \varepsilon^{-1} \int_{|x|=\varepsilon} u v \, dS \right)
\end{equation}
and
$$
\xi(u) := \xi(u,u) = \lim_{\varepsilon \to 0+} \left( \int_{|x| > \varepsilon} |\nabla u|^2 \, dx - \frac{(N-2)^2}{4} \int_{|x| > \varepsilon} \frac{u^2}{|x|^2} \, dx - \frac{N-2}{2} \varepsilon^{-1} \int_{|x|=\varepsilon} u^2 \, dS \right).
$$
Such a characterization plays a role for functions $u \in X^1 (\R^N) \setminus H^1 (\R^N)$. Otherwise, namely if $u \in H^1 (\R^N)$, it is more convenient to use the formula \eqref{normX1-oryg} instead of \eqref{normX1}.



Observe that $X^1 (\R^N)$ is not translation invariant, and that the maximal translation invariant subspace is $H^1 (\R^N)$. Indeed, if $u \in H^1 (\R^N) \subset X^1 (\R^N)$, then for every $y \in \R^N$, $u(\cdot - y) \in H^1 (\R^N) \subset X^1 (\R^N)$. Moreover we have the following fact.

\begin{Lem}
If $u \in X^1 (\R^N) \setminus H^1 (\R^N)$, then for every $y \in \R^N \setminus \{0\}$, $u(\cdot -y) \not\in X^1 (\R^N)$.
\end{Lem}

\begin{proof}
Let $u \in X^1 (\R^N) \setminus H^1 (\R^N)$. Note that, thanks to e.g. the characterization of the norm \eqref{normX1} from \cite{TrachanasZographopoulos, TrachanasZographopoulosCor}, we know that $\nabla u \in L^2 (\R^N \setminus B(0,\varepsilon))$ for every $\varepsilon > 0$. Moreover, since $u \not\in H^1 (\R^N)$, $\nabla u \not\in L^2 (B(0,\varepsilon))$ for every $\varepsilon > 0$. Fix $y \in \R^N \setminus \{0\}$ and choose $\varepsilon < \frac{|y|}{2}$. Then $\nabla u(\cdot - y) \not\in L^2 (B(y,\varepsilon))$, but $B(y,\varepsilon) \subset \R^N \setminus B(0,\varepsilon)$, so $\nabla u (\cdot -y ) \not \in L^2 (\R^N \setminus B(0,\varepsilon))$, meaning that $u (\cdot -y ) \not\in X^1 (\R^N)$.
\end{proof}

This fact makes the argument based on the Lions' concentration-compactness lemma and translations significantly more challenging than in the standard case. We overcome this difficulty by the use of continuous embeddings into $H^s (\R^N)$, which is translation-invariant, and by making use of the regularity theory to find that limits of translated Poho\v{z}aev-Palais-Smale sequences are indeed in $X^1 (\R^N)$, and even in $H^1_{\mathrm{loc}} (\R^N)$. For this purpose we show that such a limit is a solution to an elliptic problem without the singular potential in the distributional sense. Then we make use of Weyl's lemma for the Poisson problem to see that such a solution lies in $H^1_{\mathrm{loc}} (\R^N)$. Then, to show even more - the limit point belongs globally to $H^1(\R^N)$ - we make use of a variant of the Caccioppoli inequality.




\section{Variational setting}\label{sect:3}

Define $J : X^1 (\R^N) \rightarrow \R$ by
$$
J(u) = \frac12 \xi(u) - \int_{\R^N} G(u) \, dx.
$$
Then, in view of (G1)--(G3) and embeddings \eqref{eq:embeddings}, $J$ is of $\cC^1$ class and
$$
J'(u)(v) = \xi(u,v) - \int_{\R^N} g(u) v \, dx.
$$
Observe that on $X^1 (\R^N)$ we can introduce an equivalent norm
$$
\|u\|_\ell^2 := \xi(u) + \ell \int_{\R^N} u^2 \, dx
$$
and the scalar product
$$
\langle u, v \rangle_\ell := \xi(u,v) + \ell \int_{\R^N} uv \, dx.
$$
Then $J$ and its derivative can be rewritten as
$$
J(u) = \frac12 \|u\|_\ell^2 - \int_{\R^N} F(u) \, dx, \quad J'(u)(v) = \langle u,v\rangle_\ell - \int_{\R^N} f(u) v \, dx, \quad u,v \in X^1 (\R^N).
$$
Note that if $u \in X^1 (\R^N)$ is a critical point of $J$, then $u$ is a weak solution to \eqref{eq:main} on $\R^N \setminus \{0\}$. Indeed, take any test function $\varphi \in \cC_0^\infty (\R^N \setminus \{0\})$ and consider the equality $J'(u)(\varphi) = 0$. Then
$$
\xi(u, \varphi) = \int_{\R^N} g(u) \varphi \, dx.
$$
Recall \eqref{xi:uv} and, since $\supp \varphi \subset \R^N \setminus B(0,\varepsilon)$ for sufficiently small $\varepsilon > 0$, we see that
$$
\int_{|x| > \varepsilon} \nabla u \nabla \varphi \, dx = \int_{\R^N} \nabla u \nabla \varphi \, dx, \quad \int_{|x| > \varepsilon} \frac{u \varphi}{|x|^2} \, dx = \int_{\R^N} \frac{u \varphi}{|x|^2} \, dx, \quad \int_{|x|=\varepsilon} u \varphi \, dS = 0.
$$
Hence
$$
\int_{\R^N} \nabla u \nabla \varphi \, dx - \frac{(N-2)^2}{4} \int_{\R^N} \frac{u \varphi}{|x|^2} \, dx = \int_{\R^N} g(u) \varphi \, dx
$$
and $u$ is a weak solution to \eqref{eq:main} on $\R^N \setminus \{0\}$.

We are going to show that $J$ has the mountain pass geometry.

\begin{Lem}\label{r}
There is $r > 0$ such that $\inf_{\|u\|_\ell =r} J(u) > 0$. 
\end{Lem}

\begin{proof}
Note that, from (G1)--(G3) there is $C > 0$ such that $G(u) \leq - \frac{\ell}{4} |u|^2 + C |u|^{p}$. Hence
\begin{align*}
J(u) &\geq \frac12 \xi(u) + \frac{\ell}{4} \int_{\R^N} |u|^2 \, dx - C \int_{\R^N} |u|^{p} \, dx \\
&= \frac 14 \xi(u)+\frac{1}{4} \|u\|_\ell^2 - C \int_{\R^N} |u|^{p} \, dx \geq \frac{1}{4} \|u\|_\ell^2 - C \|u\|_\ell^p.
\end{align*}
Since $p > 2$, we get $\inf_{\|u\|_\ell =r} J(u) > 0$ for sufficiently small $r > 0$.
\end{proof}

\begin{Lem}\label{v}
There is $v \in X^1 (\R^N)$ such that $J(v) < 0$ and $\|v\|_\ell > r$, where $r$ is given in Lemma \ref{r}.
\end{Lem}

\begin{proof}
Following the idea from \cite{BL}, for any $R > 0$, we define
$$
w_R(x) := \left\{ \begin{array}{ll}
\zeta_0 & \quad |x| < R \\
\zeta_0 (R+1-|x|) & \quad R \leq |x| \leq R+1 \\
0 & \quad |x| > R+1
\end{array} \right.,
$$
where $\zeta_0$ is defined in (G4).
It is clear that $w_R \in H^1 (\R^N) \subset X^1 (\R^N)$ for every $R > 0$. Moreover $w_R(x) \geq 0$ for $x \in \R^N$. Recall that
$$
| B(0, R) | = c_N R^N
$$
for some $c_N > 0$. Hence $| B(0,R+1) \setminus B(0,R) | = c_N \left( (R+1)^N - R^N \right) \lesssim  R^{N-1}$. Hence
\begin{align*}
\int_{\R^N} G(w_R) \, dx &\geq G(\zeta_0) c_N R^N - \sup_{\zeta \in [0, \zeta_0]} G(\zeta) \alpha_N R^{N-1}, \quad \mbox{for some } \alpha_N > 0, \\
\int_{\R^N} |\nabla w_R|^2 \, dx & \lesssim R^{N-1}, \\
\int_{\R^N} |w_R|^2 \, dx & \gtrsim R^N.
\end{align*}
Then
$$
J(w_R) \leq \frac12 \int_{\R^N} |\nabla w_R|^2 \, dx - \int_{\R^N} G(w_R) \, dx \leq C_1 R^{N-1} - C_2 R^{N} + C_3 R^{N-1}
$$
with $C_2 > 0$. Moreover 
$$
\| w_R \|_\ell^2 \geq \ell \int_{\R^N} |w_R|^2 \, dx \gtrsim R^N.
$$
Thus $J(w_R) < 0$ and $\|w_R\|_\ell > r$ for $R > 0$ large enough.
\end{proof}

Lemma \ref{r} and Lemma \ref{v} imply that $J$ has the mountain pass geometry. We define the mountain pass level
\begin{equation}\label{def:c}
c := \inf_{\gamma \in \Gamma} \sup_{t \in [0,1]} J(\gamma(t)),
\end{equation}
where the family of paths $\Gamma$ is given by
\begin{equation}\label{def:Gamma}
\Gamma := \left\{ \gamma \in \cC ([0,1]; X^1 (\R^N)) \ : \ \gamma(0) = 0, \ J(\gamma(1)) < 0 \right\}.
\end{equation}

The following fact is a well-known consequence of the deformation lemma and we include it here for the reader's convenience.

\begin{Lem}\label{lemma:critical_point_on_path}
If there is $\gamma \in \Gamma$ and $t^*  \in (0,1)$ with 
$$
c = J(\gamma(t^*)) > J(\gamma(t)) \quad \mbox{for } t \neq t^*,
$$
then $J'(\gamma(t^*)) = 0$.
\end{Lem}

\begin{proof}
Assume by contradiction that $J'(\gamma(t^*)) \neq 0$. Since $J$ is of $\cC^1$ class, we may choose $\delta > 0$ and $\varepsilon > 0$ such that
$$
\inf \{ \| J'(v) \| \ : \ \|v - \gamma(t^*)\| \leq \delta \} > \frac{8\varepsilon}{\delta}.
$$
Then, by the deformation lemma \cite[Lemma 2.3]{W}, we may find a flow $\eta \in \cC ([0,1] \times X^1 (\R^N); X^1 (\R^N))$ such that $\eta(1, \gamma(\cdot)) \in \Gamma$ and $J(\eta(1,\gamma(t^*))) \leq c - \varepsilon < c$. Moreover, for every $t \neq t^*$, $J(\eta(1,\gamma(t))) \leq J(\gamma(t)) < c$. Hence, for every $t \in [0,1]$ we have $J(\eta(1,\gamma(t))) < c$. Thus $\sup_{t \in [0,1]} J(\eta(1,\gamma(t))) < c$ and we obtain a contradiction.
\end{proof}
We define
\[
    P(u) 
    = \xi(u) - 2^* \int_{\R^N} G(u) \, dx.
    \]

\begin{Rem}
    Note that, because of the presence of singular terms in $\xi(u)$, it is not known if all critical points $u\in X^1(\R^N)$ of $J$ satisfy the condition $P(u)=0$.
\end{Rem}




To find a Poho\v{z}aev-Palais-Smale sequence, we will use the technique of the augmented functional introduced by Jeanjean \cite{Jeanjean} (see also \cite{HirataIkomaTanaka, JeanjeanTanaka}). Then, thanks to the density of $H^1 (\R^N)$ in $X^1(\R^N)$, we can choose such a sequence in $H^1(\R^N)$, which will later simplify notations and computations, but it is not necessary for the argument to work.


\begin{Lem}\label{Lem:ExistencePPS}
There is $(u_n) \subset H^1 (\R^N)$ such that
$$
J(u_n) \to c, \quad J'(u_n) \to 0, \quad P(u_n) \to 0.
$$
\end{Lem}

\begin{proof}
Set $\Phi : \R \times X^1 (\R^N) \rightarrow X^1 (\R^N)$ given by
$$
\Phi(t, u)(x) := u(e^{-t} x).
$$
Then
$$
J(\Phi(t, u)) = \frac{e^{(N-2)t}}{2} \xi(u) - e^{Nt} \int_{\R^N} G(u) \, dx.
$$
It is clear that $J \circ \Phi$ is of class $\cC^1 (\R \times X^{1}(\R^N))$. 
Then we consider the following family of paths
$$
\widetilde{\Ga} := \left\{ \widetilde{\gamma} \in \cC ([0,1]; \R \times X^1 (\R^N)) \ : \ \widetilde{\gamma}(0)=(0,0),\, (J \circ \Phi)(\widetilde{\gamma}(1)) < 0\right\}.
$$
Note that $\Ga = \{ \Phi \circ \widetilde{\gamma} \ : \ \widetilde{\gamma} \in \widetilde{\Ga} \}$. Hence
$$
c = \inf_{\widetilde{\gamma} \in \widetilde{\Ga}} \sup_{t \in [0,1]} (J \circ \Phi)(\widetilde{\gamma}(t)).
$$
For every $\varepsilon \in \left(0, \frac{c}{2} \right)$ we may choose a path $\gamma \in \Ga$ such that $\sup_{t \in [0,1]} (J \circ \Phi)(0,\gamma(t)) \leq c + \varepsilon$. Hence, by Ekeland's variational principle \cite[Theorem 2.8]{W}, there exists $(t, u) \in \R \times X^1 (\R^N)$ such that
\begin{align*}
c - 2\varepsilon \leq (J \circ \Phi)(t,u) \leq c + 2\varepsilon, \\
\inf_{s \in [0,1]} \left( |t|^2 + \| u - \gamma(s) \|^2 \right)^{1/2} \leq 2 \sqrt{\varepsilon}, \\
\left\| (J \circ \Phi)' (t,u) \right\| \leq 2 \sqrt{\varepsilon}
\end{align*}
and we can construct a sequence $(t_n, v_n) \subset \R \times X^1 (\R^N)$ such that
$$
t_n \to 0, \quad (J \circ \Phi)(t_n, v_n) \to c, \quad (J \circ \Phi)'(t_n, v_n) \to 0.
$$
Set $\widetilde{u}_n := \Phi(t_n, v_n)$. Then, for every $(h, w) \in \R \times X^1 (\R^N)$,
$$
(J \circ \Phi)'(t_n, v_n) (h,w) = J'(\widetilde{u}_n)(\Phi(t_n,w)) + \frac{N-2}{2} P(\widetilde{u}_n)h. 
$$
It is clear that $J(\widetilde{u}_n) = J(\Phi(t_n, v_n)) \to c$. Taking $w = 0$ and $h=1$ we obtain that
$$
o(1) = (J \circ \Phi)'(t_n, v_n) (1, w) = P(\widetilde{u}_n)
$$
so that $P(\widetilde{u}_n) \to 0$. It is also classical to check that $J'(\widetilde{u}_n) \to 0$. Since $H^1(\R^N)$ is a dense subset of $X^1(\R^N)$, for each $n$ we can find $u_n$ such that
\[
\|\widetilde{u}_n-u_n\| \leq \frac 1n
\]
Then, from the continuity of $J : X^1(\R^N) \rightarrow \R$, $J' : X^1 (\R^N) \rightarrow (X^1(\R^N))^{-1}$, $P : X^1(\R^N) \rightarrow \R$, 
$$
J(u_n) \to c, \quad J'(u_n) \to 0, \quad P(u_n) \to 0.
$$
\end{proof}



\section{Profile decomposition of Poho\v{z}aev-Palais-Smale sequences}\label{sect:4}

In this section, we study the profile decomposition of bounded Poho\v{z}aev-Palais-Smale sequences. We begin with several lemmas that will be useful in proving the forthcoming decomposition result.


\begin{Lem}\label{L:scalar_product_translation}
If $u \in X^1 (\R^N)$, $w \in H^1 (\R^N)$ and $(y_n) \subset \R^N$ satisfy $|y_n| \to \infty$, then
\[
\lim_{n\to\infty} \langle u,w(\cdot-y_n)\rangle_\ell = 0.
\]
\end{Lem}

\begin{proof}
Recall that $w(\cdot -y_n) \weakto 0$ in $H^1 (\R^N)$. Indeed, fix $\varphi \in \cC_0^\infty (\R^N)$ and note that
\begin{align*}
&\quad \left| \int_{\R^N} \nabla w (\cdot - y_n) \nabla \varphi + w(\cdot - y_n) \varphi \, dx \right| = \left| \int_{-y_n + \supp \varphi} \nabla w \nabla \varphi ( (\cdot + y_n) ) + w \varphi (\cdot + y_n) \, dx \right| \\
&\leq \underbrace{\left( \int_{-y_n + \supp \varphi} |\nabla w|^2 + w^2 \, dx \right)^{1/2}}_{\to 0} \left( \int_{\R^N} |\nabla \varphi|^2 + \varphi^2 \, dx \right)^{1/2} \to 0.
\end{align*}
Now, take any $v \in H^1 (\R^N)$ and let $\varphi_k \in \cC_0^\infty (\R^N)$ be the sequence such that $\varphi_k \to v$ in $H^1(\R^N)$. Then
\begin{align*}
&\quad \left| \int_{\R^N} \nabla w (\cdot - y_n) \nabla v + w(\cdot - y_n) v \, dx \right| \\
&\leq \left| \int_{\R^N} \nabla w (\cdot - y_n) \nabla (v-\varphi_k) + w(\cdot - y_n) (v-\varphi_k) \, dx \right|  + \left| \int_{\R^N} \nabla w (\cdot - y_n) \nabla \varphi_k + w(\cdot - y_n) \varphi_k \, dx \right| \\
&\leq \|w\|_{H^1(\R^N)} \|v-\varphi_k\|_{H^1(\R^N)} + o_n(1).
\end{align*}
Now, taking $k\to\infty$ we obtain 
$$
\int_{\R^N} \nabla w (\cdot - y_n) \nabla v + w(\cdot - y_n) v \, dx \to 0.
$$
Thus, indeed $w(\cdot -y_n) \weakto 0$ in $H^1 (\R^N)$. Then, since the embedding $H^1 (\R^N) \subset X^1 (\R^N)$ is continuous, thanks to \cite[Theorem 3.10]{Brezis}, we get $w(\cdot - y_n) \weakto 0$ in $X^1 (\R^N)$. Hence $\langle u, w(\cdot - y_n) \rangle_\ell \to 0$ and the proof is completed.
\end{proof}

We will apply the following lemma several times.

\begin{Lem}[{\cite[Lemma 3.1]{GuoMederski}}] \label{L:MederskiGuo}
If $u \in H^1 (\R^N)$ and $(y_n) \subset \R^N$ satisfy $|y_n|\to\infty$, then
$$
\lim_{n\to\infty} \int_{\R^N} \frac{|u(\cdot - y_n)|^2}{|x|^2} \, dx = 0.
$$
\end{Lem}



\begin{Lem} \label{H1-translations-bounded}
If $u \in H^1 (\R^N)$ and $(y_n) \subset \R^N$ satisfy $|y_n| \to \infty$, then
$$
\lim_{n\to\infty} \| u (\cdot - y_n) \| = \| u \|_{H^1 (\R^N)}.
$$
Moreover, as a simple consequence,
$$
\lim_{n\to\infty} \|u (\cdot - y_n)\|_\ell^2 = \| \nabla u\|_{L^2 (\R^N)}^2 + \ell \| u \|_{L^2 (\R^N)}^2.
$$
\end{Lem}

\begin{proof}
Observe that 
\begin{align*}
\| u (\cdot - y_n) \|^2 &= \int_{\R^N} |\nabla u(\cdot - y_n)|^2 \, dx - \frac{(N-2)^2}{4} \int_{\R^N} \frac{|u(\cdot - y_n)|^2}{|x|^2} \, dx + \int_{\R^N} |u(\cdot - y_n)|^2 \, dx \\
&= \|u\|_{H^1 (\R^N)}^2 - \frac{(N-2)^2}{4} \int_{\R^N} \frac{|u(\cdot - y_n)|^2}{|x|^2} \, dx. 
\end{align*}
Moreover, from Lemma \ref{L:MederskiGuo},
$$
\int_{\R^N} \frac{|u(\cdot - y_n)|^2}{|x|^2} \, dx \to 0
$$
and the proof is completed.
\end{proof}

From Lions’ concentration-compactness principle, we easily obtain the following corollary, which can be viewed as a variant of the concentration–compactness principle in $X^1 (\R^N)$.

\begin{Lem}\label{lionsLem}
Let $(u_n) \subset X^1 (\R^N)$ be a bounded sequence with
$$
\lim_{n \to +\infty} \sup_{z \in \R^N} \int_{B(z, r)} |u_n|^2 \, dx = 0
$$
for some $r > 0$. Then $u_n \to 0$ in $L^t (\R^N)$ for $t \in (2, 2^*)$.
\end{Lem}

\begin{proof}
Fix $t \in (2, 2^*)$. Then there is $\varepsilon \in (0,1)$ such that $t < \frac{2N}{N-2(1-\varepsilon)}$. Since $X^1 (\R^N) \subset H^{1-\varepsilon} (\R^N)$, we know that $u_n \to 0$ in $L^s (\R^N)$ for $2 < s < \frac{2N}{N-2(1-\varepsilon)}$. Therefore, $u_n \to 0$ in $L^t (\R^N)$.
\end{proof}

In the profile decomposition presented below, we construct weak solutions to the \textit{limiting problem}
$$
-\Delta u = g(u).
$$
Since we work in $X^1 (\R^N)$, while the natural energy space for this equation is $H^1 (\R^N)$, we first need a few regularity results.

We recall the following fact, which is a consequence of Weyl's lemma.

\begin{Th}[{\cite[Lemma 3]{DiFratta}}]\label{T:dystrybucje_na_H1}
Let $\Omega \subset \R^N$ be an open set and $h \in H^{-1} (\Omega)$. Suppose that $u \in L^2 (\Omega)$ satisfies
$$
-\Delta u = h \quad \mbox{in } \cD'(\Omega).
$$
Then $u \in H^{1}_{\mathrm{loc}} (\Omega)$.
\end{Th}

Applying this theorem for $\Omega=\R^N$, we can get the following result in the setting that we are interested in.

\begin{Th}\label{th:usuwanie_osobliwosci}
Suppose that $g$ satisfies (G1)--(G3). Suppose that $u \in H^s (\R^N)$, where $s \in (0,1)$ is chosen so that $p < \frac{2N}{N-2s}$, satisfies
$$
-\Delta u = g(u) \quad \mbox{in } \cD'(\R^N ).
$$
Then $u \in H^1 (\R^N)$ and $u$ is a weak solution to $-\Delta u = g(u)$ on $\R^N$.
\end{Th}


\begin{proof}
\textbf{Step 1.} \textit{The distributional solution $u$ belongs to $H^1_{\loc} (\R^N)$}. 

Let $u$ be as in the statement. Since
\[
|g(u)| \lesssim |u|+|u|^{p-1} \in L^2(\R^N) + L^\frac{p}{p-1} (\R^N) \subset H^{-1} (\R^N),
\]
by Theorem \ref{T:dystrybucje_na_H1} we know that $u \in H^1_{\loc} (\R^N)$.

\textbf{Step 2.} \textit{The solution $u$ lies globally in $H^1 (\R^N)$.}

Now, we will show that $u \in H^1 (\R^N)$. First, we will show the following variant of the Caccioppoli inequality (cf. \cite[Chapter 4]{Giaquinta}, \cite{Iwaniec})
$$
\int_{B(0,R)} |\nabla u|^2 \, dx \lesssim \int_{B(0,2R)} |g(u) u| \, dx + \frac{1}{R^2} \int_{B(0,2R) \setminus B(0,R)} u^2 \, dx.
$$
Take $R > 0$ and $\eta_R \in \cC_0^\infty (\R^N)$ such that
$$
\eta_R \equiv 1 \mbox{ on } B(0,R), \quad \eta_R \equiv 0 \mbox{ outside of } B(0,2R)
$$
and
$$
0 \leq \eta_R \leq 1, \quad |\nabla \eta_R | \leq \frac{c}{R}
$$
for some $c \geq 1$. Then $\eta_R^2 u \in H^1_0 (B(0,2R))$. Take a sequence $\varphi_n \in \cC_0^\infty (B(0,2R))$ such that $\varphi_n \to \eta_R^2 u$ in $H^1_0(B(0,2R))$. Testing the equation with $\varphi_n$ we get
$$
\int_{\R^N} \nabla u \nabla \varphi_n \, dx = \int_{\R^N} g(u) \varphi_n \, dx.
$$
Since $\varphi_n \to \eta_R^2 u$ in the strong topology of $H^1_0(B(0,2R))$, passing to the limit, we easily get that
$$
\int_{\R^N} \nabla u \nabla (\eta_R^2 u) \, dx = \int_{\R^N} g(u) \eta_R^2 u \, dx.
$$
It can be rewritten as
$$
\int_{\R^N} \eta_R^2 |\nabla u|^2 \, dx + 2 \int_{\R^N} \eta_R u \nabla u \cdot \nabla \eta_R \, dx = \int_{\R^N} g(u) \eta_R^2 u \, dx.
$$
Now, using Young's inequality we obtain
\begin{align*}
\left| 2 \int_{\R^N} \eta_R u \nabla u \cdot \nabla \eta_R \, dx \right| &\leq  \int_{\R^N} (\eta_R |\nabla u|) (2| \nabla \eta_R| |u|) \, dx  \leq \frac{1}{2} \int_{\R^N} \eta_R^2 |\nabla u|^2 \, dx + 2 \int_{\R^N} |\nabla \eta_R|^2 |u|^2 \, dx \\
&\leq \frac12 \int_{\R^N} \eta_R^2 |\nabla u|^2 \, dx + \frac{2c^2}{R^2} \int_{B(0,2R) \setminus B(0,R)} u^2 \, dx.
\end{align*}
Hence
\begin{align*}
\int_{\R^N} \eta_R^2 |\nabla u|^2 \, dx \leq \int_{\R^N} |g(u) u| \eta_R^2 \, dx + \frac12 \int_{\R^N} \eta_R^2 |\nabla u|^2 \, dx + \frac{2c^2}{R^2} \int_{B(0,2R) \setminus B(0,R)} u^2 \, dx
\end{align*}
or equivalently
\begin{align*}
\frac12 \int_{\R^N} \eta_R^2 |\nabla u|^2 \, dx \leq \int_{B(0,2R)} |g(u) u| \eta_R^2 \, dx + \frac{2c^2}{R^2} \int_{B(0,2R) \setminus B(0,R)} u^2 \, dx.
\end{align*}
Thus
\begin{align*}
\int_{B(0,R)} |\nabla u|^2 \, dx &\leq \int_{\R^N} \eta_R^2 |\nabla u|^2 \, dx \lesssim \int_{B(0,2R)} |g(u) u| \eta_R^2 \, dx + \frac{1}{R^2} \int_{B(0,2R) \setminus B(0,R)} u^2 \, dx \\
&\leq \int_{B(0,2R)} |g(u) u| \, dx + \frac{1}{R^2} \int_{B(0,2R) \setminus B(0,R)} u^2 \, dx.
\end{align*}
Hence, we obtained the following variant of the Caccioppoli inequality 
\begin{equation}\label{caccioppoli}
    \int_{B(0,R)} |\nabla u|^2 \, dx \lesssim \int_{B(0,2R)} |g(u) u| \, dx + \frac{1}{R^2} \int_{B(0,2R) \setminus B(0,R)} u^2 \, dx.
\end{equation}
Recall that
$$
g(u) \lesssim |u| + |u|^{p-1}.
$$
Thus, from the Caccioppoli inequality \eqref{caccioppoli} we obtain
\begin{align*}
\int_{B(0,R)} |\nabla u|^2 \, dx &\lesssim \int_{B(0,2R)} |g(u) u| \, dx + \frac{1}{R^2} \int_{B(0,2R) \setminus B(0,R)} u^2 \, dx \\
&\lesssim \|u\|_{L^2 (\R^N)}^2 + \|u\|_{L^p (\R^N)}^p  + \frac{1}{R^2} \|u\|_{L^2 (\R^N)}^2
\end{align*}
Now, taking $R \to \infty$, from the monotone convergence theorem
$$
\int_{\R^N} |\nabla u|^2 \, dx \lesssim \|u\|_{L^2 (\R^N)}^2 + \|u\|_{L^p (\R^N)}^p < \infty
$$
and therefore $u \in H^1 (\R^N)$.
\end{proof}





We introduce the energy and Poho\v{z}aev functionals for the limiting equation. We define ${J_\infty : H^1 (\R^N) \rightarrow \R}$ and $P_\infty : H^1 (\R^N) \rightarrow \R$ as
\begin{align*}
J_\infty(u) &:= \frac12 \int_{\R^N} |\nabla u|^2 \, dx + \frac{\ell}{2} \int_{\R^N} u^2 \, dx - \int_{\R^N} F(u) \, dx, \\
P_{\infty}(u) &:= \frac12 \int_{\R^N} |\nabla u|^2 \, dx + \frac{2^* \ell}{2} \int_{\R^N} u^2 \, dx - 2^* \int_{\R^N}  F(u) \, dx.
\end{align*}

\begin{Rem}\label{R:problem_graniczny_wlasnosci}
It is well known that every critical point $w \in H^1 (\R^N)$ of $J_\infty$ satisfies $P_\infty(w) = 0$.
\end{Rem}





\begin{Th}\label{Th:splitting}
Suppose that $(u_n) \subset H^1 (\R^N)$ is a bounded in $X^1(\R^N)$, Poho\v{z}aev-Palais-Smale sequence for $J$ at level $c > 0$. Then, passing to a subsequence, there are an integer $k \geq 0$, $u_0 \in X^1 (\R^N)$, and sequences $(y_n^j) \subset \Z^N$, $w^j \in H^1 (\R^N)$ for $j\in \{1,\ldots, k\}$ such that

\begin{itemize}
\item[(a)] $u_n \weakto u_0$, $J'(u_0)=0$ and $P(u_0)=0$;
\item[(b)] $|y_n^j| \to \infty$ and $|y_n^j - y_n^{j'}|\to \infty$ for $j \neq j'$;
\item[(c)] $w^j \neq 0$, $w^j$ is a weak solution to
$$
-\Delta w_j = g(w_j) \quad \mbox{in } \R^N,
$$
and is a critical point of $J_{\infty}$ for each $1 \leq j \leq k$;
\item[(d)] $u_n - u_0 - \sum_{j=1}^k w^j (\cdot - y_n^j) \to 0$ in $X^1 (\R^N)$;
\item[(e)] $\|u_n\|^2_\ell\to \|u_0\|^2_\ell+\sum_{j=1}^k  \left( \|\nabla w^j\|_{L^2 (\R^N)}^2 + \ell \| w^j \|_{L^2 (\R^N)}^2 \right)$;
\item[(f)] $J(u_n) \to J(u_0) + \sum_{j=1}^k J_\infty (w^j)$.
\end{itemize}
\end{Th}




Before we present the proof, we observe that - compared with standard profile decompositions - the limiting points $w^j$ belong to the $H^1(\R^N)$ space, not merely to $X^1 (\R^N)$. This observation is crucial for the proof, since the energy functional $J_\infty$, associated with $-\Delta w = g(w)$, is defined only on $H^1 (\R^N)$ and is not well-defined for functions in $X^1 (\R^N) \setminus H^1 (\R^N)$. Without this additional regularity, we would not know that each $w^j$ is a critical point of the energy functional, and consequently, we could not ensure that the iterative procedure in the proof terminates after finitely many steps. Here, a crucial role is played by the regularity theorem, which relies on Weyl’s lemma, the Caccioppoli inequality, and continuous embeddings into $H^s (\R^N)$.


\begin{proof}


Since $(u_n) \subset X^1 (\R^N)$ is bounded, it is (up to a subsequence) weakly convergent in $X^1 (\R^N)$ to some $u_0 \in X^1 (\R^N)$. Since $J$ is weak-to-weak* continuous, for any $\varphi \in \cC_0^\infty (\R^N)$, $J'(u_0)(\varphi) = J'(u_n)(\varphi) + o(1) = o(1)$, so $J'(u_0) = 0$.



Let $v_n^1 := u_n - u_0 \weakto 0$ in $X^1 (\R^N)$. Then $v_n^1 \to 0$ in $L^t_{loc} (\R^N)$ for $t \in [2,2^*)$ and $v_n^1 (x) \to 0$ for a.e. $x \in \R^N$. It is clear that
\begin{equation}\label{E:v_n^1-norm}
    \|v_n^1\|^2_\ell = \| u_n - u_0\|^2_\ell = \|u_n\|^2_\ell + \|u_0\|^2_\ell - 2 \langle u_n, u_0 \rangle_\ell = \|u_n\|^2_\ell - \|u_0\|^2_\ell + o (1).
\end{equation}

Now, we compute
\begin{align*}
\int_{\R^N} F(u_n) - F(v_n^1) \, dx  &= \int_{\R^N} \int_0^1 f(u_n - s u_0)u_0 \, ds \, dx = \int_0^1 \int_{\R^N} f(u_n - s u_0)u_0 \, dx \, ds.
\end{align*}
The family $\{ f(u_n - s u_0)u_0 \}$ is uniformly integrable and tight on $\R^N$, so that from the Vitali convergence theorem
\begin{equation}\label{E:F_convergence}
\begin{split}
\int_{\R^N} F(u_n) - F(v_n^1) \, dx &\to
 \int_0^1 \int_{\R^N} f(u_0 - s u_0)u_0 \, dx \, ds \\
&=   \int_{\R^N} \int_0^1 f(u_0 - s u_0)u_0  \, ds \, dx = \int_{\R^N} F(u_0) \, dx.
\end{split}
\end{equation}
Hence
$$
J(v_n^1) = J(u_n) - J(u_0) + o(1).
$$
Arguing similarly we obtain that $P(v_n^1) = P(u_n)-P(u_0)+o(1)=-P(u_0)+o(1)$.

\textbf{Step 1.}
Suppose that there is $r > 0$ such that
\begin{equation}\label{lionsCond}
\lim_{n \to +\infty} \sup_{z \in \R^N} \int_{B(z, r)} |v_n^1|^2 \, dx = 0.
\end{equation}
From Lemma \ref{lionsLem} there follows that $v_n^1 \to 0$ in $L^t (\R^N)$ for $t \in (2,2^*)$.
From (F1)--(F3), for every $\varepsilon > 0$ we have
$$
|f(u)| \leq \varepsilon |u| + C_\varepsilon |u|^{p-1}.
$$
Then

\begin{align*}
o(1) &= J'(u_n)(v_n^1) = \langle u_n, u_n-u_0 \rangle_\ell - \int_{\R^N} f(u_n)(u_n-u_0) \, dx \\
&= \| v_n^1 \|^2_\ell + \langle u_0, u_n - u_0 \rangle_\ell - \int_{\R^N} f(u_n)(u_n-u_0) \, dx.
\end{align*}
Clearly $\langle u_0, u_n - u_0 \rangle_\ell \to 0$ and from the H\"older inequality
\begin{align*}
\limsup_{n\to \infty} &\left| \int_{\R^N} f(u_n^+)(u_n-u_0) \, dx \right| \leq \limsup_{n\to \infty} \left( \varepsilon \int_{\R^N} |u_n^+| |u_n-u_0| \, dx +  C_\varepsilon \int_{\R^N} |u_n^+|^{p-1} |u_n-u_0| \, dx \right) \\ 
&\leq \limsup_{n\to \infty} \left( \varepsilon \| u_n^+ \|_{L^2 (\R^N)} \|u_n-u_0\|_{L^2 (\R^N)} + C_\varepsilon \| u_n^+ \|_{L^p (\R^N)}^{p-1} \|u_n-u_0\|_{L^p (\R^N)} \right) \\
&= \varepsilon \limsup_{n\to\infty} \| u_n^+ \|_{L^2 (\R^N)} \|u_n-u_0\|_{L^2 (\R^N)}.
\end{align*}
Taking $\varepsilon \to 0^+$ and using the boundedness of $(u_n)$ in $L^2 (\R^N)$ we get that
$$
\int_{\R^N} f(u_n)(u_n-u_0) \, dx \to 0.
$$
Thus $v_n^1\to 0$ i.e. $u_n \to u_0$ in $X^1 (\R^N)$ and the proof is completed with $k = 0$. Indeed, conditions (b) and (c) are empty, (d) and (e) were just proven with empty sums, (f) and $P(u_0)=0$ follow from the continuity of $P$.


\textbf{Step 2.}
Suppose that \eqref{lionsCond} does not hold. Then we may find $(y_n^1) \subset \Z^N$ and $n_0$ such that
$$
\inf_{n \geq n_0} \int_{B(y_n^1, r+\sqrt{N})} |v_n^1|^2 \, dx > 0
$$ 
It is standard to verify that $|y_n^1| \to \infty$. Set $w_n^1 := u_n (\cdot + y_n^1)\geq 0$. It is clear that $(w_n^1) \subset H^1 (\R^N)$, since $H^1 (\R^N)$ is a translation-invariant subspace. 


Note that $0<\frac{N(p-2)}{2(p-1)}<1$. Therefore, we can choose $s\in \left( \max \left\{ \frac{N(p-2)}{2(p-1)}, \frac12 \right\}, 1 \right)$. Then we have $p<\frac{2N-2}{N-2}<\frac{2N}{N-2s}$. Then $(u_n)$ is bounded in $H^s (\R^N)$, and since $H^s (\R^N)$ is translation-invariant, $(w_n^1) \subset H^s (\R^N)$ is also bounded in $H^s (\R^N)$. Hence, there is $w^1 \in H^s (\R^N)$ such that $w_n^1 \weakto w^1$ in $H^s (\R^N)$, $w_n^1 \to w^1$ in $L^t_{loc} (\R^N)$ for $t \in [2,2^*_s)$ and $w_n^1 (x) \to w^1 (x)$ for a.e. $x \in \R^N$. 


Observe that
$$
\inf_n \int_{B(0, r+\sqrt{N})} |w_n^1|^2 \, dx 
=\inf_n \int_{B(y_n^1, r+\sqrt{N})} |u_n|^2 \, dx
\geq \inf_n \int_{B(y_n^1, r+\sqrt{N})} |v_n^1|^2+|u_0|^2 \, dx
> 0,
$$
and therefore $w^1 \neq 0$. Then, from Lemma \ref{H1-translations-bounded}, for every test function $\varphi \in \cC_0^\infty (\R^N)$
$$
|J'(u_n) (\varphi(\cdot - y_n^1))| \leq \| \nabla J(u_n) \|_\ell \| \varphi (\cdot - y_n^1) \|_\ell \to 0 \cdot \left( \| \nabla \varphi \|_{L^2 (\R^N)}^2 + \ell \| \varphi \|_{L^2 (\R^N)}^2 \right)^{1/2} = 0,
$$
and
\begin{align*}
&\quad o(1) = J'(u_n) (\varphi(\cdot - y_n^1)) \\
&= \int_{\R^N} \nabla u_n \nabla \varphi(\cdot - y_n^1) \, dx - \frac{(N-2)^2}{4} \int_{\R^N} \frac{u_n \varphi(\cdot - y_n^1)}{|x|^2} \, dx - \int_{\R^N} g(u_n) \varphi(\cdot - y_n^1) \, dx = (*).
\end{align*}
Hence, fix a test function $\varphi \in \cC_0^\infty (\R^N)$ and choose $n_0$ such that for $n \geq n_0$ we have\linebreak $\mathrm{dist}\, (\supp \varphi, -y^1_n) > 1$. Then, for $n \geq n_0$ we also have
$$
\int_{\R^N} \frac{u_n \varphi(\cdot - y_n^1)}{|x|^2} \, dx = \int_{\R^N} \frac{u_n(\cdot + y_n^1) \varphi}{|x+y_n^1|^2} \, dx = \int_{\supp \varphi} \frac{u_n(\cdot + y_n^1) \varphi}{|x+y_n^1|^2} \, dx.
$$
Similarly
$$
\int_{\R^N} \nabla u_n \nabla \varphi(\cdot - y_n^1) \, dx = \int_{\R^N} \nabla u_n(\cdot + y_n^1) \nabla \varphi \, dx = \int_{\supp \varphi} \nabla u_n(\cdot + y_n^1) \nabla \varphi \, dx.
$$
Thus
\begin{align*}
\mathclap{
(*) = \int_{\supp \varphi} \nabla u_n(\cdot + y_n^1) \nabla \varphi \, dx - \frac{(N-2)^2}{4} \int_{\supp \varphi} \frac{u_n(\cdot + y_n^1) \varphi}{|x+y_n^1|^2} \, dx - \int_{\supp \varphi} g( u_n(\cdot + y_n^1) ) \varphi \, dx = (**).}
\end{align*}
From Lebesgue dominated convergence theorem we get
$$
\int_{\supp \varphi} g( u_n(\cdot + y_n^1) ) \varphi \, dx \to \int_{\supp \varphi} g( w^1) \varphi \, dx.
$$
Since $\supp \varphi$ is compact and $|y_n^1| \to \infty$, for every $n$ we can choose $m(n)$ such that $m(n) \to \infty$ and $|x+y_n^1| \geq m(n)$ for $x\in\supp \varphi$. Then
$$
\int_{\supp \varphi} \frac{|u_n(\cdot + y_n^1) \varphi|}{|x+y_n^1|^2} \, dx \leq \frac{1}{m(n)^2} \int_{\supp \varphi} |u_n(\cdot + y_n^1) \varphi| \, dx \leq \frac{1}{m(n)^2} \|u_n\|_{L^2 (\R^N)} \| \varphi\|_{L^2 (\R^N)} \to 0,
$$
since $(u_n)$ is bounded in $L^2 (\R^N)$. Then
\begin{align*}
(**) = \int_{\supp \varphi} \nabla u_n(\cdot + y_n^1) \nabla \varphi \, dx - \int_{\supp \varphi} g( w^1) \varphi \, dx + o(1).
\end{align*}
Hence, we get
$$
\int_{\R^N} \nabla w_n^1 \nabla \varphi \, dx \to \int_{\R^N} g(w^1) \varphi \, dx
$$
for all $\varphi \in \cC_0^\infty (\R^N)$. We can rewrite the left hand side as
$$
\int_{\R^N} \nabla w_n^1 \nabla \varphi \, dx = - \int_{\R^N} w_n^1 \Delta \varphi \, dx
$$
and therefore, since $w_n^1 \to w^1$ in $L^2_{\mathrm{loc}}(\R^N)$, we obtain
$$
- \int_{\R^N} w^1 \Delta \varphi \, dx = \int_{\R^N} g(w^1) \varphi \, dx,
$$
namely $w^1 \in H^s (\R^N)$ is a solution to
$$
-\Delta w^1 = g(w^1) \quad \mbox{in } \cD' (\R^N).
$$
From Theorem \ref{T:dystrybucje_na_H1}, we obtain that $w^1 \in H^1 (\R^N)$. In particular, $w^1$ is a weak solution to the problem
\begin{equation}\label{eq:limiting}
-\Delta w^1 = g(w^1)
\end{equation}
in $\R^N$. Therefore, by Remark \ref{R:problem_graniczny_wlasnosci}, $J_\infty '(w^1) = 0$ and $P_\infty (w^1) = 0$.



\textbf{Step 3.} Put $v_n^2=u_n-u_0-w^1(\cdot - y_n^1)=v_n^1-w^1(\cdot - y_n^1)$. Note that, by Lemma \ref{H1-translations-bounded} and \eqref{E:v_n^1-norm} we have
\begin{equation*}
\begin{split}
\|v_n^2\|^2_{\ell}&=\|v_n^1\|^2_\ell+\|w^1(\cdot-y_n^1)\|^2_\ell - 2\langle v_n^1,w^1(\cdot-y_n^1)\rangle_\ell
\\
&=\|u_n\|^2_\ell-\|u_0\|^2_\ell+o(1)+\|w^1(\cdot -y_n^1)\|^2_\ell
-2\langle u_n,w^1(\cdot -y_n^1)\rangle_\ell
+2\langle u_0,w^1(\cdot -y_n^1)\rangle_\ell\\
&=\|u_n\|^2_\ell-\|u_0\|^2_\ell+\left( \|\nabla w^1\|_{L^2 (\R^N)}^2 + \ell \| w^1 \|_{L^2 (\R^N)}^2 \right) - 2\langle u_n,w^1(\cdot -y_n^1)\rangle_\ell +o(1)=(*),
\end{split}
\end{equation*}
where the term $2\langle u_0,w^1 (\cdot - y_n^1)\rangle_\ell $ disappears by Lemma \ref{L:scalar_product_translation}. Then, we compute
\begin{align*}
\langle u_n, w^1 (\cdot - y_n^1) \rangle_\ell = &\int_{\R^N} \nabla u_n(\cdot + y_n^1) \nabla w^1 \,dx + \ell \int_{\R^N}  u_n(\cdot + y_n^1)  w^1 \,dx \\ &- \frac{(N-2)^2}{4} \int_{\R^N} \frac{u_n w^1 (\cdot - y_n^1)}{|x|^2} \, dx.
\end{align*}
We will show that
\begin{equation}\label{E:temp-new}
\int_{\R^N} \frac{u_n w^1 (\cdot - y_n^1)}{|x|^2} \, dx\to 0
\end{equation}
to finally get
\[
\begin{split}
(*)=&\|u_n\|^2_\ell-\|u_0\|^2_\ell+\left( \|\nabla w^1\|_{L^2 (\R^N)}^2 + \ell \| w^1 \|_{L^2 (\R^N)}^2 \right)\\
&-2 \left( \int_{\R^N} \nabla u_n(\cdot + y_n^1) \nabla w^1 \,dx + \ell \int_{\R^N}  u_n(\cdot + y_n^1)  w^1 \,dx  \right) +o(1) = (**).
\end{split}
\]
Now, using the fact that $w^1$ is a solution to \eqref{eq:limiting}, we obtain
\begin{align}\label{E:temp-new2}
\int_{\R^N} \nabla u_n(\cdot + y_n^1) \nabla w^1 \,dx + \ell \int_{\R^N}  u_n(\cdot + y_n^1)  w^1 \,dx = \int_{\R^N} f(w^1) w_n^1 \, dx.
\end{align}
Since $w_n^1$ is bounded in $L^2 (\R^N)$ and in $L^p (\R^N)$, $w_n^1 \to w^1$ in $L^2_{\mathrm{loc}} (\R^N)$ and in $L^p_{\mathrm{loc}} (\R^N)$, from Vitali convergence theorem
$$
\int_{\R^N} f(w^1) w_n^1 \, dx \to \int_{\R^N} f(w^1) w^1 \, dx = \|\nabla w^1\|_{L^2 (\R^N)}^2 + \ell \| w^1 \|_{L^2 (\R^N)}^2.
$$
Thus
\begin{align*}
(**) &= \|u_n\|^2_\ell-\|u_0\|^2_\ell - \left( \|\nabla w^1\|_{L^2 (\R^N)}^2 + \ell \| w^1 \|_{L^2 (\R^N)}^2 \right) + o(1).
\end{align*}
To prove \eqref{E:temp-new} take any $R>0$ and consider
\[
\int_{\R^N} \frac{u_n w^1 (\cdot - y_n^1)}{|x|^2} \, dx
=\int_{B(0,R)} \frac{u_n w^1 (\cdot - y_n^1)}{|x|^2} \, dx
+\int_{\R^N\setminus B(0,R)} \frac{u_n w^1 (\cdot - y_n^1)}{|x|^2} \, dx=: I_1+I_2
\]
By H\"{o}lder inequality
\[
|I_1|\leq \|u_n\|_{L^{\frac{2N}{N-2s}}(\R^n)} \cdot 
\left(\int_{B(0,R)} \frac{|w^1(x-y_n)|^{\frac{2N}{N+2s}}}{|x|^{\frac{4N}{N+2s}}} \, dx\right)^{\frac{N+2s}{2N}}.
\]
By the continuity of the embedding $X^1(\R^N)\subset H^s(\R^N)$, the sequence $(u_n)$ is bounded in $L^{\frac{2N}{N-2s}} (\R^N)$. 

Since $w^1 \in H^1 (\R^N)$ is a weak solution to \eqref{eq:limiting}, from \cite[Theorem 2]{PankovDecay}, there is $\alpha > 0$ such that $|w^1(x)| \lesssim e^{-\alpha |x|}$. For $x\in B(0,R)$ we have $|x-y_n|\geq |y_n|-|x|\geq |y_n|-R\to\infty$, so one has
     \[
     \delta_n = \sup_{x\in B(0,R)} |w^1(x-y_n)| \lesssim  e^{-\alpha (|y_n|-R))} \to 0 \quad \mbox{as } n\to\infty.
     \]
Moreover, since $s > \frac12$,
$$
\frac{4N}{N+2s} < N.
$$
Thus $\int_{B(0,R)} |x|^{- \frac{4N}{N+2s}} \, dx < +\infty$. Hence
$$
|I_1| \lesssim \delta_n \to 0 \quad \mbox{as } n \to \infty.
$$
Since $u_n, w^1\in L^2 (\R^N)$, and $(u_n)$ is bounded in $L^2 (\R^N)$, we have
\[
|I_2| \leq \frac{1}{R^2} \|u_n\|_{L^2(\R^N)}\|w^1\|_{L^2(\R^N)}.
\]
Thus, for any $R > 0$, we get
$$
\limsup_{n\to\infty} \left| \int_{\R^N} \frac{u_n w^1 (\cdot - y_n^1)}{|x|^2} \, dx \right| \lesssim \frac{1}{R^2}.
$$
Hence,
$$
\int_{\R^N} \frac{u_n w^1 (\cdot - y_n^1)}{|x|^2} \, dx \to 0
$$
and the proof of \eqref{E:temp-new} is completed.



Arguing as in \eqref{E:F_convergence} for the pair of functions $(v_n^1, w^1)$ instead of $(u_n, u_0)$, and then directly from \eqref{E:F_convergence} we get
\begin{equation}\label{E:nonlinearity_decomposition}
\begin{split}
\int_{\R^N} F(v_n^2) \, dx 
&=\int_{\R^N} F(v_n^2(\cdot + y_n^1)) \, dx \\
&=\int_{\R^N} F(v_n^1(\cdot+y_n^1))\,dx -\int_{\R^N} F(w_1)\,dx +o(1)\\
&=\int_{\R^N} F(v_n^1)\,dx -\int_{\R^N} F(w_1)\,dx +o(1)\\
&=\int_{\R^N} F(u_n)\,dx -\int_{\R^N} F(u_0)\,dx - \int_{\R^N} F(w_1)\,dx +o(1).
\end{split}
\end{equation}
Combining the foregoing with $\|v_n^2\|_\ell=\|u_n\|_\ell^2 - \|u_0\|_\ell^2 - \left( \|\nabla w^1\|_{L^2 (\R^N)}^2 + \ell \| w^1 \|_{L^2 (\R^N)}^2 \right) +o(1)$ we obtain
\[
J(v_n^2)=J(u_n)-J(u_0)-J_\infty(w_1)+o(1)
\]
and, by Remark \ref{R:problem_graniczny_wlasnosci},
\[
P(v_n^2)=P(u_n)-P(u_0)-P_\infty(w_1)+o(1)=-P(u_0)+o(1).
\]
Consider 
\[
\Theta_2=\lim_{n \to +\infty} \sup_{z \in \R^N} \int_{B(z, r)} |v_n^2|^2 \, dx.
\]
If $\Theta_2=0$, similarly as in \eqref{lionsCond} and Step 1, using $J'(u_n)(v_n^2)=o(1)$, \eqref{E:temp-new}, and \eqref{E:temp-new2}, we get $v_n^2\to 0$ in $X^1(\R^N)$ and the proof is completed with $k=1$.

Suppose $\Theta_2>0$. Then we reason as in Step 2. We find the sequence $(y_n^2)\subset \Z^N$, $|y_n^2|\to\infty$ such that 
$$
\inf_{n \geq n_0} \int_{B(y_n^2, r+\sqrt{N})} |v_n^2|^2 \, dx > 0.
$$ 
It is classical to check that $|y_n^1-y_n^2|\to\infty$. We put $w_n^2=u_n(\cdot+y_n^2)\in H^1(\R^N)$. Then we find $w^2\in H^1(\R^N)$ such that 
\begin{itemize}
\item $w_n^2\weakto w^2$ in $H^s(\R^N)$ for $s$ sufficiently close to $1$,
\item $w_n^2\to w^2$ in $L^t_{loc}(\R^N)$ for $t\in[2,2^*_s)$,
\item $w^2$ is a weak solution of $-\Delta w^2 =g(w^2)$ on $\R^N$,
\item $J'_\infty(w^2)=0, \ P_\infty(w^2)=0$.
\end{itemize}
Put 
\[
v_n^3:=v_n^2-w^2(\cdot - y_n^2)
=u_n-u_0-w^1(\cdot - y_n^1)-w^2(\cdot - y_n^2).
\]
The same way as in \eqref{E:nonlinearity_decomposition} we get
\[
\int_{\R^N} F(v_n^3) \, dx 
=\int_{\R^N} F(u_n)\,dx 
-\int_{\R^N} F(u_0)\,dx 
- \int_{\R^N} F(w_1)\,dx 
- \int_{\R^N} F(w_2)\,dx +o(1).
\]
Moreover, by Lemma \ref{L:scalar_product_translation}, Lemma \ref{H1-translations-bounded}, and \eqref{E:temp-new}
\[
\begin{split}
\|v_n^3\|^2_\ell=&\|v_n^2\|^2_\ell + \|w^2(\cdot - y_n^2)\|^2_\ell
-2\langle v_n^2, w^2(\cdot - y_n^2)\rangle_\ell\\
=& \|v_n^2\|^2_\ell 
+\left( \|\nabla w^2\|_{L^2 (\R^N)}^2 + \ell \| w^2 \|_{L^2 (\R^N)}^2 \right)\\
&-2\left( \langle u_n,w^2(\cdot - y_n^2)\rangle_\ell
+\langle u_0,w^2(\cdot - y_n^2)\rangle_\ell\right)
-2\langle w^1(\cdot - y_n^1),w^2(\cdot - y_n^2)\rangle_\ell
\\
=&\|v_n^2\|^2_\ell 
+\left( \|\nabla w^2\|_{L^2 (\R^N)}^2 + \ell \| w^2 \|_{L^2 (\R^N)}^2 \right)\\
&-2\left(\langle u_n(\cdot + y_n^2), w^1 \rangle_\ell
+\frac{(N-2)^2}4\int_{\R^N} \frac{|w^2|^2 }{|x|^2} \, dx\right)
-2\langle w^1(\cdot - y_n^1),w^2(\cdot - y_n^2)\rangle_\ell
+o(1)\\
=&\|v_n^2\|^2_\ell 
-\left( \|\nabla w^2\|_{L^2 (\R^N)}^2 + \ell \| w^2 \|_{L^2 (\R^N)}^2 \right)-2\langle w^1(\cdot - y_n^1),w^2(\cdot - y_n^2)\rangle_\ell+o(1).
\end{split}
\]
By Cauchy-Schwarz inequality and Lemma \ref{L:MederskiGuo} we have
\[
\left| \int_{\R^N} \frac{ w^1(\cdot - y_n^1)w^2(\cdot - y_n^2)}{|x|^2} \, dx \right| \leq \left(
\int_{\R^N} \frac{ (w^1(\cdot - y_n^1))^2}{|x|^2} \, dx \right)^{1/2} \left(
\int_{\R^N} \frac{ (w^2(\cdot - y_n^1))^2}{|x|^2} \, dx \right)^{1/2}\to 0
\]
and 
\[
\left|\int_{\R^N} \frac{ w^1w^2(\cdot - (y_n^2-y_n^2))}{|x|^2} \, dx \right| \leq \left(
\int_{\R^N} \frac{ (w^1)^2}{|x|^2} \, dx \right)^{1/2} \left(
\int_{\R^N} \frac{ (w^2(\cdot - (y_n^1-y_n^2)))^2}{|x|^2} \, dx \right)^{1/2} \to 0.
\]
Therefore, again by Lemma $\ref{L:MederskiGuo}$ and utilizing $|y_n^1-y_n^2|\to\infty$ we get
\[
\langle w^1(\cdot - y_n^1),w^2(\cdot - y_n^2)\rangle_\ell=\langle w^1,w^2(\cdot - (y_n^2-y_n^1))\rangle_\ell + o(1)\to 0.
\]
Hence,
\[
\begin{split}
\|v_n^3\|^2_\ell&=\|v_n^2\|^2_\ell 
-\left( \|\nabla w^2\|_{L^2 (\R^N)}^2 + \ell \| w^2 \|_{L^2 (\R^N)}^2 \right)+o(1)\\
&=\|u_n\|^2_\ell-\|u_0\|^2_\ell-
\sum_{j=1}^2 \left( \|\nabla w^j\|_{L^2 (\R^N)}^2 + \ell \| w^j \|_{L^2 (\R^N)}^2 \right)+o(1).
\end{split}
\]
Finally we can write
\[
J(v_n^3)=J(u_n)-J(u_0)-\sum_{j=1}^2 J_\infty(w^j)+o(1)
\]
and
\[
P(v_n)=P(u_n)-P(u_0)-\sum_{j=1}^2 P_\infty(w^j)=P(u_0)+o(1).
\]

\textbf{Step 4.} We can repeat the above procedure to construct, for some $k\in\Z$:
\begin{itemize}
\item the sequences $(y_n^j)\subset \Z^N$, $j=1,\ldots,k$ such that $|y_n^j|\to\infty$ and $|y_n^j-y_n^{j'}|\to \infty$ for $j\neq j'$,
\item $w_j\in H^1(\R^N)$, $j=1,\ldots,k$, such that $P_\infty(w^j)=0$,
\item $\|u_n\|^2_\ell\to \|u_0\|^2_\ell+\sum_{j=1}^k  \left( \|\nabla w^j\|_{L^2 (\R^N)}^2 + \ell \| w^j \|_{L^2 (\R^N)}^2 \right)$.
\end{itemize}
Note that, from (F1)--(F3), for every $\varepsilon > 0$, there is a constant $C_\varepsilon > 0$ such that
\begin{align*}
\| \nabla w^j \|_{L^2(\R^N)}^2 &+ \ell \| w^j \|_{L^2(\R^N)}^2 = \int_{\R^N} f(w^j)w^j \, dx \leq \varepsilon \|w^j\|_{L^2(\R^N)}^2 + C_\varepsilon \|w^j\|_{L^p(\R^N)}^p \\
&\leq C \varepsilon \left( \| \nabla w^j \|_{L^2(\R^N)}^2 + \ell \| w^j \|_{L^2(\R^N)}^2 \right) + C C_\varepsilon \left( \| \nabla w^j \|_{L^2(\R^N)}^2 + \ell \| w^j \|_{L^2(\R^N)}^2 \right)^{p/2},
\end{align*}
where we used that $J_\infty '(w^j)(w^j) = 0$ and $C > 0$ comes from Sobolev embeddings. Hence, choosing sufficiently small $\varepsilon > 0$,
$$
\frac{1 - C\varepsilon}{C C_\varepsilon} \leq \left(\| \nabla w^j \|_{L^2 (\R^N)}^2 + \ell \| w^j \|_{L^2(\R^N)}^2 \right)^{(p-2)/2}.
$$
It means that for every $j$,
$$
\| \nabla w^j \|_{L^2(\R^N)}^2 + \ell \| w^j \|_{L^2(\R^N)}^2 \gtrsim 1,
$$
hence the procedure needs to terminate after a finite number of steps, since $(u_n)$ is bounded in $X^1(\R^N)$.

\end{proof}

\section{Existence of a solution}\label{sect:5}

We will start by showing that the Poho\v{z}aev-Palais-Smale sequences are bounded. 

\begin{Lem}\label{Lem:PPS-bounded}
Suppose that $c > 0$ and $(u_n) \subset H^1 (\R^N)$ is a sequence such that $J(u_n) \to c$, $J'(u_n) \to 0$ and $P(u_n) \to 0$. Then $(u_n)$ is bounded in $X^1 (\R^N)$.
\end{Lem}

\begin{proof}
Since $J(u_n) \to c$, and using $P(u_n) \to 0$, we get
$$
\frac12 \xi (u_n) =  \int_{\R^N} G(u_n) \, dx + c + o(1) = \frac{1}{2^*} \xi(u_n) + c + o(1)
$$
and therefore
$$
\left(\frac12 - \frac{1}{2^*}\right) \xi(u_n) = c + o(1). 
$$
In particular, $(\xi(u_n))$ is bounded. Now, what remains is to show that $\|u_n\|_{L^2 (\R^N)}$ is bounded. For this purpose, using (G5), we consider

$$
J(u_n) - \frac{1}{\nu} J'(u_n)(u_n) = \left( \frac{1}{2} - \frac{1}{\nu} \right) \xi(u_n) - \int_{\R^N} \left( G(u_n) - \frac{1}{\nu} g(u_n) u_n \right) \, dx \geq \left( \frac{1}{2} - \frac{1}{\nu} \right) \xi(u_n) - \gamma_0 \|u_n\|_2^2.
$$
Hence
\begin{align}\label{ineq1}
- \gamma_0 \|u_n\|_{L^2(\R^N)}^2 \leq J(u_n) - \frac{1}{\nu} J'(u_n)(u_n) - \left( \frac{1}{2} - \frac{1}{\nu} \right) \xi(u_n) \leq C - \frac{1}{\nu} J'(u_n)(u_n)
\end{align}
for some $C > 0$. Since $J'(u_n) \to 0$, we get that 
\begin{align*}
\left| \frac{1}{\nu} J'(u_n)(u_n) \right| \leq \delta_n \|u_n\|_\ell = \delta_n \left( \xi(u_n) + \|u_n\|_{L^2(\R^N)}^2 \right)^{1/2} &\leq \delta_n \left( C + \|u_n\|_{L^2(\R^N)}^2 \right)^{1/2} \\
&\leq \delta_n \left(C^{1/2} + \|u_n\|_{L^2(\R^N)}\right).
\end{align*}
for some $C>0$ and $\delta_n \to 0$. Finally, combining it with \eqref{ineq1} we find that
$$
-\gamma_0 \|u_n\|_{L^2(\R^N)}^2 \leq C + C \|u_n\|_{L^2 (\R^N)}
$$
and $\|u_n\|_{L^2(\R^N)}$ is bounded, since $\gamma_0 < 0$.
\end{proof}

Now we are ready to prove the existence of a nontrivial solution.

\begin{proof}[Proof of Theorem \ref{Th:Main1-Existence}(a)]

Define $c_\infty := \inf_{\gamma \in \Gamma_\infty} \sup_{t \in [0,1]} J_\infty(\gamma(t))$, where
$$
\Gamma_\infty := \left\{ \gamma \in C([0,1]; H^1 (\R^N)) \ : \ \gamma(0) = 0, \ J_\infty(\gamma(1)) < 0 \right\}.
$$
It is known that $c_\infty > 0$ and that there is an optimal path $\gamma_\infty \in \Gamma_\infty$ such that
$$
c_\infty = \sup_{t\in [0,1]} J_\infty (\gamma_\infty(t)),
$$
see \cite{JeanjeanTanaka}. Moreover
$$
c_\infty = \inf \left\{ J_\infty(w) \ : \ w \in H^1 (\R^N) \setminus \{0\}, \ J_\infty'(w) = 0 \right\}.
$$
Since $J(u) < J_\infty(u)$ for $u \in H^1 (\R^N) \setminus \{0\}$ and $\Gamma_\infty \subset \Gamma$, we obtain the following
\begin{align*}
c_\infty = \sup_{t\in [0,1]} J_\infty (\gamma_\infty(t)) > \sup_{t\in [0,1]} J (\gamma_\infty(t)) \geq \inf_{\gamma \in \Gamma} \sup_{t\in[0,1]} J (\gamma(t)) = c,
\end{align*}
where $c$ and $\Gamma$ are given by \eqref{def:c} and \eqref{def:Gamma}, respectively. 



Let $(u_n)$ be a Poho\v{z}aev-Palais-Smale sequence obtained in Lemma \ref{Lem:ExistencePPS}. From Lemma \ref{Lem:PPS-bounded} this sequence is bounded in $X^1 (\R^N)$ and, up to choosing a subsequence, (a)--(f) of Theorem \ref{Th:splitting} hold. Since $w^j \in H^1 (\R^N)$ are nontrivial critical points of $J_\infty$, we know that $J_\infty(w^k) \geq c_\infty$. Hence, from (f),
$$
J(u_n) = J(u_0) + \sum_{j=1}^k J_\infty (w^j) + o(1) \geq J(u_0) + k c_\infty + o(1).
$$
Hence 
$$
c \geq J(u_0) + k c_\infty > J(u_0) + k c.
$$
Since $u_0$ is a critical point of $J$ satisfying $P(u_0) = 0$, $J(u_0) = \left( \frac12 - \frac{1}{2^*} \right) \xi(u_0) \geq 0$, and therefore $k = 0$. Then, from (d) and (e), $u_n \to u_0$ in $X^1 (\R^N)$, $J(u_0) = c$ and $u_0$ is a nontrivial solution to \eqref{eq:main}. Moreover, the solution $u_0$ satisfies $P(u_0) = 0$.
\end{proof}


\section{Further properties of a solution}\label{sect:6}

\subsection{\texorpdfstring{$u_0$}{u0} is a minimizer on Poho\v{z}aev constraint}

Recall the Poho\v{z}aev constraint
$$
\cM := \left\{ u \in X^1 (\R^N) \setminus \{0\} \ : \ \xi(u) = 2^* \int_{\R^N} G(u) \, dx \right\}.
$$
We will show that the solution $u_0$ is a minimizer on $\cM$, namely that there holds
$$
J(u_0) = \inf_{\cM} J.
$$

\begin{proof}[Proof of Theorem \ref{Th:Main1-Existence}(b)]
Since $u_0 \in \cM$, clearly $c = J(u_0) \geq m := \inf_{\cM} J$. Hence it is enough to show that $m \geq c$. For that purpose, take any $u \in \cM$ and consider the rescaled function $U(\theta) := u(\cdot / \theta)$ for $\theta \geq 0$, where we set $U(0) = 0$. Then
$$
\| U(\theta) \|_\ell^2 = \theta^{N-2} \xi(u) + \theta^N \ell \|u\|_{L^2(\R^N)}^2
$$
and
$$
\Psi(\theta) := J(U(\theta)) = \frac{\theta^{N-2}}{2} \xi(u) - \theta^N \int_{\R^N} G(u) \, dx.
$$
Then, clearly $U : [0,\infty) \rightarrow X^1 (\R^N)$ is continuous. Since $u \in \cM$, we get
$$
\int_{\R^N} G(u) \, dx = \frac{1}{2^*} \xi(u) > 0.
$$
The last inequality holds, because there is no $L^2(\R^N)$ function that optimizes the Hardy inequality. Now, we consider $\frac{d}{d \theta} \Psi$. Note that
$$
\frac{d}{d \theta} \Psi(\theta) = \frac{(N-2) \theta^{N-3}}{2 } \xi(u) - N \theta^{N-1} \int_{\R^N} G(u) \, dx = \theta^{N-3} \left( \frac{N-2}{2 } \xi(u) - N \theta^{2} \int_{\R^N} G(u) \, dx \right).
$$
Since $u \in \cM$, we get $\frac{d}{d \theta} \Psi(1) = 0$, and clearly $\frac{d}{d \theta} \Psi(\theta) > 0$ for $\theta \in (0,1)$ and $\frac{d}{d \theta} \Psi(\theta) < 0$ for $\theta \in (1,\infty)$.

Hence, the following properties are satisfied by $U$,
$$
U(0) = 0, \ J(U(\vartheta)) < 0, \ U(1) = u \in U([0,\vartheta])
$$
for sufficiently large $\vartheta > 1$; moreover $\max_{\theta \in [0,\vartheta]} J(U(\theta)) = J(u)$. In particular, the path 
\begin{equation}\label{eq:path}
\gamma := U(\vartheta \cdot)
\end{equation}
belongs to $\Gamma$, where $\Gamma$ is given by \eqref{def:Gamma}. Hence $c \leq J(u)$. Since $u \in \cM$ was arbitrary, $c \leq m$ and the proof is completed.
\end{proof}

\subsection{Regularity of the solution}

\begin{proof}[Proof of Theorem \ref{Th:Main1-Existence}(c)]
Let $u_0\in X^1(\R^N)$ be a weak solution of \eqref{eq:main}. Fix $x_0\in\R^N\setminus\{0\}$ and its open neighborhood $\Omega\subset\overline{\Omega}\subset \R^N \setminus \{0\}$. Testing $J'(u_0)=0$ on the function $\varphi\in C_0^\infty(\Omega)$ we easily get that $u_0\in H^1(\Omega)$ is a weak solution of
\begin{equation}\label{E:C_2_temp}
    -\Delta u=g(u)+\frac{(N-2)^2}{4|x|^2} u \quad\text{on } \Omega.
\end{equation}

Observe that 
\[
\left|g(u)+\frac{(N-2)^2}{4|x|^2} u \right|\lesssim |u|+|u|^{p-1} \leq (1+|u|^{p-2})(1+|u|)
\]
and $|u|^{p-2}\in L^{N/2}(\Omega)$. Therefore, by Brezis-Kato lemma (see \cite[Lemma B.3]{Struwe}), $u\in L^q_{\loc}(\Omega)$ for any $q<\infty$. Then, by Cald\'{e}ron-Zygmund inequality (see \cite[Lemma B.2]{Struwe}), $u\in W^{2,q}_{\loc}(\Omega)$ and next, by Sobolev embedding theorem, $u\in C^{1,\alpha}_{\loc}(\Omega)$. Since $x_0$ was an arbitrary point, we conclude $u\in C^{1,\alpha}_{\loc}(\R^N\setminus\{0\})$.
\end{proof}

\begin{proof}[Proof of Theorem \ref{Th:Main1-Existence}(d)]
We have already shown that $u\in C^{1,\alpha}_{\loc}(\R^N \setminus \{0\})$. Since $g$ is H\"older continuous, by Schauder estimates we get $u\in C^{2}(\R^N \setminus \{0\})$.
\end{proof}


\subsection{There is a radial solution, which is a minimizer on \texorpdfstring{$\cM$}{M}}

For a nonnegative, measurable function $u$, let $u^*$ denote its radially decreasing rearrangement (also known as the Schwarz symmetrization). The well-known Pólya-Szegő inequality states that if $u \in H^1 (\R^N)$, then $u^* \in H^1 (\R^N)$ and
$$
\int_{\R^N} |\nabla u^*|^2 \, dx \leq \int_{\R^N} |\nabla u|^2 \, dx.
$$
Moreover, the following properties hold
$$
\int_{\R^N} \frac{|u^*|^2}{|x|^2} \, dx \geq \int_{\R^N} \frac{u^2}{|x|^2} \, dx, \quad \int_{\R^N} G(u^*) \, dx = \int_{\R^N} G(u) \, dx.
$$
Observe that these inequalities imply that
$$
\xi(u^*) \leq \xi(u), \quad J(u^*) \leq J(u)
$$
as long as $u \in H^1 (\R^N)$. We will show that it also extends to the space $X^1 (\R^N)$.

\begin{Th}[Pólya-Szegő inequality in $X^1 (\R^N)$]
For any $u \in X^1 (\R^N)$, $u \geq 0$, there holds
$$
\xi(u^*) \leq \xi(u).
$$
\end{Th}

\begin{proof}
Take any $u \in X^1 (\R^N) \setminus H^1 (\R^N)$, $u \geq 0$ and take a sequence $u_n \in H^1 (\R^N)$, $u_n \geq 0$ such that $u_n \to u$ in $X^1 (\R^N)$. Consider $v_n := (u_n)^*$. Then, clearly
$$
\xi(v_n) \leq \xi(u_n), \quad \| v_n \| \leq \| u_n \|.
$$
Since $u_n \to u$ in $X^1(\R^N)$, $(v_n)$ is bounded in $X^1 (\R^N)$ and therefore, up to a subsequence, $v_n \weakto w$ in $X^1 (\R^N)$ and $v_n \to w$ in $L^2_{\mathrm{loc}}(\R^N)$ for some $w \in X^1 (\R^N)$. Since $v_n \to u^*$ in $L^2 (\R^N)$, from the uniqueness of the limit in $L^2 (\R^N)$, $w = u^*$ almost everywhere. Hence $u^* \in X^1 (\R^N)$. Since $\| \cdot \|$ is weakly lower semicontinuous,
$$
\| u^* \|^2 \leq \liminf_{n\to\infty} \|v_n\|^2 \leq \liminf_{n\to\infty} \|u_n\|^2 = \|u\|^2.  
$$
Taking into account that $\| u^* \|_{L^2 (\R^N)} = \| u \|_{L^2 (\R^N)}$, we obtain that
$$
\xi(u^*) \leq \xi(u).
$$
\end{proof}


\begin{proof}[Proof of Theorem \ref{Th:Main1-Existence}(e)]
We assume in addition that $g$ is odd. Let $u_0 \in X^1 (\R^N)$ be the solution that minimizes $J$ on $\cM$, and let $\gamma_0$ denote the path given by \eqref{eq:path} associated with $u_0$. Then we consider the Schwarz symmetrization $|\gamma (t)|^*$ of $|\gamma(t)|$ for every $t \in [0,1]$. Then
$$
|\gamma (0)|^* = 0^* = 0, \ J(|\gamma(1)|^*) \leq J(|\gamma(1)|) \leq J(\gamma(1)) < 0
$$
and therefore $|\gamma (\cdot )|^* \in \Gamma$, where $\Gamma$ is given by \eqref{def:Gamma}. Hence
$$
c \leq \max_{t \in [0,1]} J (|\gamma (t)|^*).
$$
On the other hand
$$
\max_{t \in [0,1]} J (|\gamma (t)|^*) \leq \max_{t \in [0,1]} J (|\gamma (t)|)\leq \max_{t \in [0,1]} J(\gamma(t)) = J(u_0) = c.
$$
and
$$
J (|\gamma (t)|^*) \leq J(|\gamma(t)|) \leq J(\gamma(t)) < c
$$
for $t \neq \frac{1}{\vartheta}$, where $\vartheta$ is the same as in \eqref{eq:path}. Thus
$$
J (\gamma^* (1/\vartheta) ) = c.
$$
Define $v_0 := \gamma^* (1/\vartheta)$. Then, from Lemma \ref{lemma:critical_point_on_path}, $v_0$ is a critical point of $J$ and $J(v_0) = c$. To see that $v_0 \in \cM$ observe that
$$
v_0 = \gamma^*(1/\vartheta) = |u_0|^*
$$
and from $J(v_0) = J(u_0)$ we immediately obtain $\xi(v_0) = \xi(u_0)$. Hence
$$
\xi(v_0) = \xi(u_0) = 2^* \int_{\R^N} G(u_0) \, dx = 2^* \int_{\R^N} G(v_0) \, dx
$$
and $v_0 \in \cM$. From Theorem \ref{Th:Main1-Existence}(c), $v_0 \in C^{1,\alpha}_{\loc}(\R^N \setminus \{0\})$ for some $\alpha \in (0,1)$. Since $v_0$ is a radial function (with a small abuse of notation, $v_0(x)=v_0(r)$, $r=|x|$) we can rewrite \eqref{eq:main}, in a weak sense, as
\[
-v_0''-\frac{N-1}{r} v_0'=g(v_0)+\frac{(N-2)^2}{4r^2} v_0 \
\Longleftrightarrow \ -\left(v_0'r^{N-1}\right)'=r^{N-1}\left(g(v_0)+\frac{(N-2)^2}{4r^2} v_0\right).
\]
The right-hand side is continuous in $r$, hence by the standard bootstrap argument the map $r \mapsto v_0'(r)r^{N-1}\in C^1$ on its domain, and therefore $v_0\in C^2$ in a neighborhood of $x_0$. Since $x_0$ was an arbitrary point, we conclude $v_0\in C^2(\R^N\setminus\{0\})$.
\end{proof}





\subsection{Nonnegative solutions do not belong to \texorpdfstring{$H^1(\R^N)$}{H1}}

\begin{proof}[{Proof of Theorem \ref{Th:Main1-Existence}(f)}]
Let $u \in X^1 (\R^N) \cap C^2 (\R^N \setminus \{0\})$ be a nonnegative solution to \eqref{eq:main}. As in the argument in the proof of \cite[Theorem 1.1(iv)]{LiLiTang}, we introduce
$$
u(x) = |x|^\tau w(x),
$$
where $\tau\in\left(-\frac{N-2}2,0\right)$. Then, $w$ solves
$$
- \div \left( |x|^{2\tau} \nabla w \right) = |x|^\tau \left( \frac{\left(\tau + \frac{N-2}{2} \right)^2}{|x|^2} |x|^\tau w + g(|x|^\tau w) \right).
$$
Since, by (G1)--(G3), $g(s) \gtrsim -s$ for $s \geq 0$, we find a radius $\rho > 0$ such that
$$
\frac{\left(\tau + \frac{N-2}{2} \right)^2}{|x|^2} |x|^\tau w + g(|x|^\tau w)  \geq 0 \quad \mbox{for} \ 0 < |x| < \rho. 
$$
Then
$$
- \div \left( |x|^{2\tau} \nabla w \right) \geq 0 \quad \mbox{in } B(0,\rho) \setminus \{0\}.
$$
Repeating the argument in \cite{LiLiTang}, 
$$
u(x) = |x|^\tau w(x) \gtrsim |x|^\tau \quad \mbox{in a neighborhood of the origin},
$$
i.e. $u(x)\to \infty$ for $x\to 0$. Therefore, we can choose a radius $\widetilde{\rho} > 0$ such that $u(x) > 0$ and $g(u(x)) \geq 0$ in $B(0,\widetilde{\rho}) \setminus \{0\}$. Hence, $u$ satisfies the inequality
$$
-\Delta u - \frac{(N-2)^2}{4 |x|^2} u \geq 0 \quad \mbox{in } B(0,\widetilde{\rho}) \setminus \{0\}.
$$
From \cite[Proposition 2.1]{Smets}, $u \not\in H^1 (B(0,\widetilde{\rho}))$.
\end{proof}



\section{Existence of non-radial solutions}\label{sect:7}

In this section we will prove Theorem \ref{Th:Main2}. Following \cite{MR4173560}, we recall the subspace $X_\tau$ given by \eqref{eq:Xtau} and the group $\mathcal{O} = \mathcal{O}(M) \times \mathcal{O}(M) \times \mathcal{O}(N-2M)$ acting on $\R^N = \R^M \times \R^M \times \R^{N-2M}$, $N \geq 4$. Repeating the reasoning from \cite[Corollary 3.2]{MR4173560} we see that the following concentration-compactness principle holds in the subspace $H^s_\mathcal{O} (\R^N)$ of $\cO$-invariant functions in $H^s (\R^N)$.

\begin{Cor}
Suppose that $(u_n) \subset H^s_\mathcal{O} (\R^N)$ is bounded and for all $r > 0$
$$
\lim_{n\to\infty} \sup_{z \in \R^{N-2M}} \int_{B((0,0,z),r)} |u_n|^2 \, dx = 0.
$$
Then $u_n \to 0$ in $L^p (\R^N)$ for any $p \in \left(2, \frac{2N}{N-2s} \right)$.
\end{Cor}
Hence, with the same reasoning as in the proof of Lemma \ref{lionsLem}, we obtain the following concentration-compactness principle in $X^1_\mathcal{O} (\R^N)$.
\begin{Cor}\label{lionsLemWithSymmetry}
Suppose that $(u_n) \subset X^1_\mathcal{O} (\R^N)$ is bounded and for all $r > 0$
$$
\lim_{n\to\infty} \sup_{z \in \R^{N-2M}} \int_{B((0,0,z),r)} |u_n|^2 \, dx = 0.
$$
Then $u_n \to 0$ in $L^p (\R^N)$ for any $p \in \left(2, 2^* \right)$.
\end{Cor}

Now, we can repeat the proof of Theorem \ref{Th:splitting}, using Corollary \ref{lionsLemWithSymmetry} instead of Lemma \ref{lionsLem}, and obtain the following theorem.

\begin{Th}\label{th:dekompNonradial}
Suppose that $(u_n) \subset X_\tau \cap H^1_{\mathcal{O}} (\R^N)$ is a bounded in $X^1(\R^N)$, Poho\v{z}aev-Palais-Smale sequence for $J$ at level $c > 0$. Then, passing to a subsequence, there are an integer $k \geq 0$, $u_0 \in X_\tau \cap X^1_{\mathcal{O}} (\R^N)$, $u_0\geq 0$, and sequences $(y_n^j) \subset \R^{N-2M}$, $w^j \in X_\tau \cap H^1_{\mathcal{O}} (\R^N)$ for $j\in \{1,\ldots, k\}$ such that

\begin{itemize}
\item[(a)] $u_n \weakto u_0$, $J'(u_0)=0$ and $P(u_0)=0$;
\item[(b)] $|y_n^j| \to \infty$ and $|y_n^j - y_n^{j'}|\to \infty$ for $j \neq j'$;
\item[(c)] $w^j \neq 0$ and $w^j$ is a weak solution to
$$
-\Delta w_j = g(w_j) \quad \mbox{in } \R^N
$$
and is a critical point of $J_{\infty}$ for each $1 \leq j \leq k$;
\item[(d)] $u_n - u_0 - \sum_{j=1}^k w^j (\cdot - (0,0,y_n^j)) \to 0$ in $X^1 (\R^N)$;
\item[(e)] $\|u_n\|^2_\ell\to \|u_0\|^2_\ell+\sum_{j=1}^k  \left( \|\nabla w^j\|_{L^2 (\R^N)}^2 + \ell \| w^j \|_{L^2 (\R^N)}^2 \right)$;
\item[(f)] $J(u_n) \to J(u_0) + \sum_{j=1}^k J_\infty (w^j)$.
\end{itemize}
\end{Th}

The slight difference in the proof is that, we know $w^j$ are weak solutions to the limiting problem in $X_\tau \cap H^1_\mathcal{O} (\R^N)$, so a critical point of $J_\infty |_{X_\tau \cap H^1_{\mathcal{O}}}$. However, from Palais' principle of symmetric criticality, it is a critical point of the free functional $J_\infty$. Similarly we know that $u_0$ is a critical point of $J$.

In the same way as in Lemma \ref{r}, there is $\rho > 0$ such that
$$
\inf_{\|u\|_\ell = \rho} J(u) > 0.
$$
To show that $J$ has a mountain pass geometry, in Lemma \ref{v} it is shown how to, for any $R > 0$, construct a radial function $w_R \in H^1(B(0,R+1)) \cap L^\infty (\R^N) \subset X^1 (\R^N)$ such that 
$$
\int_{\R^N} G(w_R) \, dx \geq C_1 R^N - C_2 R^{N-1}, \quad C_1, C_2 > 0,
$$
and
$$
\|w_R\|_{L^2 (\R^N)}^2 \gtrsim R^N.
$$
Let $R > 1$. Now, as in \cite[Remark 4.2]{MR4173560}, we choose an odd and smooth function $\varphi : \R \rightarrow [-1,1]$ such that $\varphi(x) = 1$ for $x \geq 1$, and $\varphi(x) = -1$ for $x \leq -1$. Then, as shown in \cite[Remark 4.2]{MR4173560}, the function (with a small abuse of notation $w_R (x) = w_R( |x| )$) 
$$
\widetilde{w}_R (x_1, x_2, x_3) := w_R ( \sqrt{|x_1|^2+ |x_2|^2+ |x_3|^2}) \varphi(|x_1|-|x_2|)
$$
belongs to $X_\tau \cap H^1_{\mathcal{O}} (\R^N)$ and
$$
\int_{\R^N} G(\widetilde{w}_R) \, dx \geq \int_{\R^N} G(w_R) \, dx - C_3 \sum_{i=N-M}^{N-1} R^i \geq C_1 R^N - C_2 \sum_{i=N-M}^{N-1} R^i.
$$
We compute
\begin{align*}
\int_{\R^N} | \widetilde{w}_R |^2 \, dx =& \int_{\R^N} |w_R|^2 | \varphi (|x_1|-|x_2|) |^2 \, dx \\
%=& \int_0^\infty \int_0^\infty \int_0^\infty | w_R (r) |^2 | \varphi (r_1-r_2) |^2 r_1^{M-1} r_2^{M-1} r_3^{N-2M-1} \, dr_1 \, dr_2 \, dr_3 \\
=& \int_0^{R+1} \int_0^{R+1} \int_0^{R+1} | w_R (r) |^2 | \varphi (r_1-r_2) |^2 r_1^{M-1} r_2^{M-1} r_3^{N-2M-1} \, dr_1 \, dr_2 \, dr_3 \\
=& 2 \int_0^{R+1} \int_0^{R+1} \int_{r_2}^{R+1} | w_R (r) |^2 | \varphi (r_1-r_2) |^2 r_1^{M-1} r_2^{M-1} r_3^{N-2M-1} \, dr_1 \, dr_2 \, dr_3 \\
% =& \underbrace{2 \int_0^{R+1} \int_0^{R+1} \int_{r_2}^{R+1} | w_R (r) |^2 r_1^{M-1} r_2^{M-1} r_3^{N-2M-1} \, dr_1 \, dr_2 \, dr_3}_{= \int_{\R^N} |w_R|^2 \, dx} \\
% &- 2 \int_0^{R+1} \int_0^{R+1} \int_{r_2}^{r_2+1} | w_R (r) |^2  r_1^{M-1} r_2^{M-1} r_3^{N-2M-1} \, dr_1 \, dr_2 \, dr_3 \\
% &+ 2 \int_0^{R+1} \int_0^{R+1} \int_{r_2}^{r_2+1} | w_R (r) |^2 | \varphi (r_1-r_2) |^2 r_1^{M-1} r_2^{M-1} r_3^{N-2M-1} \, dr_1 \, dr_2 \, dr_3 \\
% &= \int_{\R^N} |w_R|^2 \, dx - 2 \int_0^{R+1} \int_0^{R+1} \int_{r_2}^{r_2+1} | w_R (r) |^2  r_1^{M-1} r_2^{M-1} r_3^{N-2M-1} \, dr_1 \, dr_2 \, dr_3 \\
% &+ 2 \int_0^{R+1} \int_0^{R+1} \int_{r_2}^{r_2+1} | w_R (r) |^2 | \varphi (r_1-r_2) |^2 r_1^{M-1} r_2^{M-1} r_3^{N-2M-1} \, dr_1 \, dr_2 \, dr_3 \\
=& \int_{\R^N} |w_R|^2 \, dx \\
&- 2 \int_0^{R+1} \int_0^{R+1} \int_{r_2}^{r_2+1} \underbrace{| w_R (r) |^2 (1 - | \varphi (r_1-r_2) |^2)}_{ \leq \zeta_0^2 } r_1^{M-1} r_2^{M-1} r_3^{N-2M-1} \, dr_1 \, dr_2 \, dr_3 \\
% &\geq \int_{\R^N} |w_R|^2 \, dx - 2 \int_0^{R+1} \int_0^{R+1} \int_{r_2}^{r_2+1} \underbrace{| w_R (r) |^2}_{\leq \zeta_0^2}  r_1^{M-1} r_2^{M-1} r_3^{N-2M-1} \, dr_1 \, dr_2 \, dr_3 \\
\geq& \int_{\R^N} |w_R|^2 \, dx - 2 \zeta_0^2 \int_0^{R+1} \int_0^{R+1} \int_{r_2}^{r_2+1}  r_1^{M-1} r_2^{M-1} r_3^{N-2M-1} \, dr_1 \, dr_2 \, dr_3 \\ 
=& \int_{\R^N} |w_R|^2 \, dx - \sum_{i=N-M}^{N-1} \beta_i R^i
\end{align*}
for some $\beta_i > 0$. Hence, for sufficiently large $R > 1$ we have 
$$
J(\widetilde{w}_R) < 0, \quad \|\widetilde{w}_R\|_\ell^2 > \rho.
$$

\begin{proof}[Proof of Theorem \ref{Th:Main2}]

From the above reasoning, $J$ has the mountain pass geometry, and repeating the same arguments as in Lemma \ref{Lem:ExistencePPS} we find a Poho\v{z}aev-Palais-Smale sequence for $J$ on the mountain pass level. Then, using Theorem \ref{th:dekompNonradial} repeating the proof of Theorem \ref{Th:Main1-Existence}(a), we find a nontrivial critical point of $J |_{X_\tau \cap X^1_{\mathcal{O}} (\R^N)}$, which - thanks to Palais' principle of symmetric criticality - is also a critical point of $J$. Moreover the solution belongs to $\cM$, and - as in Theorem \ref{Th:Main1-Existence}(b) - it is a minimizer of $J$ on $X_\tau \cap H^1_{\mathcal{O}} (\R^N) \cap \cM$.

It remains to show
$$
\inf_{\cM \cap X_\tau \cap X^1_{\mathcal{O}} (\R^N)} J \geq 2 \inf_\cM J.
$$

Let $u \in \cM \cap X_\tau \cap X^1_{\mathcal{O}} (\R^N)$ be the minimizer found above. Define two open sets
$$
\Omega_1 := \{ x \in \R^N \ : \ |x_1| > |x_2| \}, \quad \Omega_2 := \{ x \in \R^N \ : \ |x_1| < |x_2| \}.
$$
Define $u_1 := \mathbbm{1}_{\Omega_1} u$, $u_2 := \mathbbm{1}_{\Omega_2} u$. Then $u_1, u_2 \in X^1 (\R^N)$, since $u = 0$ on $\{|x_1|=|x_2|\}$, and $\xi(u) = 2 \xi(u_1) = 2 \xi(u_2)$. We also have $\int_{\R^N} G(u) \, dx = 2 \int_{\R^N} G(u_1) \, dx$. Hence $u_1 \in \cM$, $J(u_1) = J(u_2)$ and since $u_1$, $u_2$ have disjoint supports, we get
$$
J(u) = J(u_1) + J(u_2) = 2 J(u_1) \geq 2 \inf_{\cM} J.
$$
\end{proof}


\section*{Acknowledgements}
Bartosz Bieganowski and Daniel Strzelecki were partly supported by the National Science Centre, Poland (Grant No. 2022/47/D/ST1/00487).
\bibliographystyle{acm}
\bibliography{Bibliography}

\end{document}
